% Generated by the blueprint tool (blueprint-pdf/1) from a blueprint snapshot; edit the
% Markdown sources, not this file.  Compile with `latexmk -pdf` or by running
% `pdflatex` twice (the second run settles the contents and the page references).
% Images, if any, are read from assets/ next to this file.
\documentclass[11pt]{article}
\usepackage[T1]{fontenc}
\usepackage[utf8]{inputenc}
\usepackage{lmodern}
\usepackage{textcomp}
\usepackage[a4paper,margin=2.5cm,headheight=14pt]{geometry}
\usepackage{amsmath,amssymb,amsthm}
\usepackage{amscd}  % KaTeX's CD diagrams; part of amsmath, so always there
\usepackage{mathtools}
\usepackage{bm}
\usepackage{graphicx}
\usepackage{xcolor}
\usepackage{array,booktabs,tabularx}
\usepackage{fancyvrb}
\usepackage[normalem]{ulem}
\usepackage{microtype}
% Optional, and only with their fonts: a small TeX installation has the
% packages without the fonts, which fails at the first formula.
\IfFileExists{mathrsfs.sty}{\IfFileExists{rsfs10.tfm}{\usepackage{mathrsfs}}{}}{}
\IfFileExists{stmaryrd.sty}{\IfFileExists{stmary10.tfm}{\usepackage{stmaryrd}}{}}{}
\IfFileExists{cancel.sty}{\usepackage{cancel}}{}
\usepackage[hyphens]{url}
\usepackage[pdfusetitle,bookmarksnumbered=false]{hyperref}
\usepackage{bookmark}
\definecolor{bplink}{HTML}{1F4E9C}
\definecolor{bpmuted}{HTML}{5F676E}
\hypersetup{colorlinks=true,linkcolor=bplink,urlcolor=bplink,citecolor=bplink}
\setlength{\parindent}{0pt}
\setlength{\parskip}{0.55em plus 0.2em minus 0.1em}
% Lists as tight as the page's.
\makeatletter
\def\@listi{\leftmargin\leftmargini\parsep 0pt\topsep 0.3em plus 0.1em\itemsep 0.1em plus 0.1em}
\let\@listI\@listi
\def\@listii{\leftmargin\leftmarginii\labelwidth\leftmarginii\advance\labelwidth-\labelsep\parsep 0pt\topsep 0.1em\itemsep 0.1em}
% An ordered list that does not start at 1.
\newcommand{\bpstartat}[1]{\setcounter{\@enumctr}{\numexpr#1-1\relax}}
\makeatother
\setcounter{tocdepth}{3}
\emergencystretch=3em
\allowdisplaybreaks

% A numbered statement, set as a paper sets one: "Theorem 2.1 (Title)."
% #1 is the lead ("Theorem~2.1"), #2 what follows it in roman.
\newenvironment{bpstatement}[2]{\par\addvspace{0.6\baselineskip}\noindent\textbf{#1}#2\textbf{.}\enspace\ignorespaces}
  {\par\addvspace{0.4\baselineskip}}
\newenvironment{bpstatementit}[2]{\begin{bpstatement}{#1}{#2}\itshape\ignorespaces}{\end{bpstatement}}
% A step: a relation with prose of its own, read under the object it starts from.
\newenvironment{bpstep}[1]{\par\addvspace{0.3\baselineskip}\begin{list}{}{\leftmargin=1.6em\rightmargin=0pt\topsep=0pt}\item\relax
  {\small\color{bpmuted}#1}\par\nopagebreak}{\end{list}\par}
% Notes about where an object is written out.
\newenvironment{bpnote}{\par\small\color{bpmuted}\itshape}{\par}
\newcommand{\bpomitted}[1]{\par{\small\color{bpmuted}\itshape #1\par}}
\newcommand{\bpbodyheading}[1]{\par\addvspace{0.4\baselineskip}\noindent\textbf{#1}\par\nopagebreak}
\newcommand{\bprule}{\par\medskip\noindent\hfil\rule{0.3\linewidth}{0.4pt}\hfil\par\medskip}
\newcommand{\bpboundary}[1]{\par\nopagebreak{\small\color{bpmuted}#1\par}\nopagebreak}
% A reference to something this export leaves out, and one that resolves nowhere.
\newcommand{\bpout}[1]{#1\textsuperscript{\hyperlink{bp-outside}{\ensuremath{\ast}}}}
\newcommand{\bpbroken}[1]{\textcolor{red!70!black}{[#1?]}}
\newcommand{\bpslug}[1]{\texttt{#1}}
% A bibliography entry: label, anchor, entry; the label hangs in the margin.
\newcommand{\bpbibitem}[3]{\par\setbox0\hbox{[#1]\enspace}\ifdim\wd0<5em\setbox0\hbox to 5em{[#1]\hfil}\fi
  \hangindent=\wd0 \hangafter=1 \noindent\box0\phantomsection\label{#2}#3\par\smallskip}
% A character the export knows no way to set.
\newcommand{\bpglyph}[1]{\fbox{\scriptsize #1}}
% Images, at their own size unless they are wider than the line.
\newsavebox{\bpimagebox}
\newcommand{\bpimage}[1]{\sbox{\bpimagebox}{\includegraphics{#1}}%
  \ifdim\wd\bpimagebox>\linewidth\resizebox{\linewidth}{!}{\usebox{\bpimagebox}}\else\usebox{\bpimagebox}\fi}
\newcommand{\bpmissingimage}[1]{\fbox{\parbox{0.8\linewidth}{\centering\small\color{bpmuted}Image not included: #1}}}
\DefineVerbatimEnvironment{bpverbatim}{Verbatim}{fontsize=\small,xleftmargin=1em}
\renewcommand{\qedsymbol}{\ensuremath{\blacksquare}}
\hypersetup{pdftitle={Weil Conjectures — 14 Purity of real sheaves},pdfsubject={Chapter 14 Purity of real sheaves; all levels},pdfcreator={blueprint}}
% KaTeX commands the website knows
\providecommand{\R}{\mathbb{R}}\providecommand{\N}{\mathbb{N}}\providecommand{\Z}{\mathbb{Z}}
\providecommand{\Q}{\mathbb{Q}}\providecommand{\C}{\mathbb{C}}
\providecommand{\Reals}{\mathbb{R}}\providecommand{\reals}{\mathbb{R}}\providecommand{\Complex}{\mathbb{C}}
\providecommand{\cnums}{\mathbb{C}}\providecommand{\natnums}{\mathbb{N}}\providecommand{\Bbb}[1]{\mathbb{#1}}
\providecommand{\bold}[1]{\mathbf{#1}}\providecommand{\frak}[1]{\mathfrak{#1}}\providecommand{\mathbbm}[1]{\mathbb{#1}}
\providecommand{\lang}{\langle}\providecommand{\rang}{\rangle}\providecommand{\bra}[1]{\langle #1\rvert}
\providecommand{\ket}[1]{\lvert #1\rangle}\providecommand{\braket}[1]{\langle #1\rangle}
\providecommand{\set}[1]{\{#1\}}\providecommand{\Set}[1]{\left\{#1\right\}}
\providecommand{\infin}{\infty}\providecommand{\isin}{\in}\providecommand{\plusmn}{\pm}\providecommand{\sub}{\subset}
\providecommand{\sube}{\subseteq}\providecommand{\supe}{\supseteq}\providecommand{\exist}{\exists}
\providecommand{\bull}{\bullet}\providecommand{\larr}{\leftarrow}\providecommand{\rarr}{\rightarrow}
\providecommand{\lrarr}{\leftrightarrow}\providecommand{\Larr}{\Leftarrow}\providecommand{\Rarr}{\Rightarrow}
\providecommand{\Lrarr}{\Leftrightarrow}\providecommand{\uarr}{\uparrow}\providecommand{\darr}{\downarrow}
\providecommand{\harr}{\leftrightarrow}\providecommand{\lArr}{\Leftarrow}\providecommand{\rArr}{\Rightarrow}
\providecommand{\hArr}{\Leftrightarrow}\providecommand{\alef}{\aleph}\providecommand{\alefsym}{\aleph}
\providecommand{\real}{\Re}\providecommand{\image}{\Im}\providecommand{\weierp}{\wp}
\providecommand{\Alpha}{\mathrm{A}}\providecommand{\Beta}{\mathrm{B}}\providecommand{\Epsilon}{\mathrm{E}}
\providecommand{\Zeta}{\mathrm{Z}}\providecommand{\Eta}{\mathrm{H}}\providecommand{\Iota}{\mathrm{I}}
\providecommand{\Kappa}{\mathrm{K}}\providecommand{\Mu}{\mathrm{M}}\providecommand{\Nu}{\mathrm{N}}
\providecommand{\Omicron}{\mathrm{O}}\providecommand{\Rho}{\mathrm{P}}\providecommand{\Tau}{\mathrm{T}}
\providecommand{\Chi}{\mathrm{X}}\providecommand{\omicron}{o}\providecommand{\gt}{>}\providecommand{\lt}{<}
\providecommand{\And}{\&}\providecommand{\hdashline}{\hline}\providecommand{\operatornamewithlimits}[1]{\operatorname*{#1}}
\providecommand{\llbracket}{[\![}\providecommand{\rrbracket}{]\!]}\providecommand{\cancel}[1]{#1}
\providecommand{\mathscr}{\mathcal}
% the characters beyond ASCII this document uses
\DeclareUnicodeCharacter{2160}{\textup{I}}
\DeclareUnicodeCharacter{2161}{\textup{II}}
\DeclareUnicodeCharacter{2162}{\textup{III}}
% the project's macros ([katex.macros] in blueprint.toml)
\providecommand{\A}{}\renewcommand{\A}{\mathbf A}
\providecommand{\AEn}{}\renewcommand{\AEn}{A_{E,n}}
\providecommand{\Aaf}{}\renewcommand{\Aaf}{\mathbf A}
\providecommand{\Ab}{}\renewcommand{\Ab}{\mathsf{Ab}}
\providecommand{\AffSch}{}\renewcommand{\AffSch}{\operatorname{AffSch}}
\providecommand{\Aone}{}\renewcommand{\Aone}{\mathbf A^1}
\providecommand{\Bl}{}\renewcommand{\Bl}{\operatorname{Bl}}
\providecommand{\Br}{}\renewcommand{\Br}{\mathrm{Br}}
\providecommand{\C}{}\renewcommand{\C}{\mathbb C}
\providecommand{\Cech}{}\renewcommand{\Cech}{\check{H}}
\providecommand{\Coext}{}\renewcommand{\Coext}{\operatorname{Coext}}
\providecommand{\Cone}{}\renewcommand{\Cone}{\operatorname{Cone}}
\providecommand{\Cons}{}\renewcommand{\Cons}{\operatorname{Cons}}
\providecommand{\Cref}{}\renewcommand{\Cref}[1]{\text{#1}}
\providecommand{\D}{}\renewcommand{\D}{D}
\providecommand{\DLAFT}{}\renewcommand{\DLAFT}{\operatorname{FT}}
\providecommand{\DLAFr}{}\renewcommand{\DLAFr}{\operatorname{Fr}}
\providecommand{\DLAQ}{}\renewcommand{\DLAQ}{{\overline{\mathbf Q}_\ell}}
\providecommand{\DLRQ}{}\renewcommand{\DLRQ}{{\overline{\mathbf Q}_\ell}}
\providecommand{\Dbc}{}\renewcommand{\Dbc}{D^b_c}
\providecommand{\Dc}{}\renewcommand{\Dc}{D_{\mathrm{cons}}}
\providecommand{\Dctf}{}\renewcommand{\Dctf}{D^b_{ctf}}
\providecommand{\Dic}{}\renewcommand{\Dic}{D_{\IndCons}}
\providecommand{\Dplus}{}\renewcommand{\Dplus}{D^{+}}
\providecommand{\End}{}\renewcommand{\End}{\operatorname{End}}
\providecommand{\Et}{}\renewcommand{\Et}{\mathrm{\acute{e}t}}
\providecommand{\Ex}{}\renewcommand{\Ex}{\mathrm{Ex}}
\providecommand{\Ext}{}\renewcommand{\Ext}{\operatorname{Ext}}
\providecommand{\F}{}\renewcommand{\F}{\mathbb F}
\providecommand{\FT}{}\renewcommand{\FT}{\operatorname{FT}}
\providecommand{\Fbar}{}\renewcommand{\Fbar}{{\overline{\mathbf F}_q}}
\providecommand{\Fell}{}\renewcommand{\Fell}{\mathbf F_\ell}
\providecommand{\Fib}{}\renewcommand{\Fib}{\operatorname{Fib}}
\providecommand{\Fix}{}\renewcommand{\Fix}{\operatorname{Fix}}
\providecommand{\Fp}{}\renewcommand{\Fp}{{\mathbf F_{p}}}
\providecommand{\Fq}{}\renewcommand{\Fq}{{\mathbf F_q}}
\providecommand{\Fqn}{}\renewcommand{\Fqn}{{\mathbf F_{q^{n}}}}
\providecommand{\Frob}{}\renewcommand{\Frob}{\operatorname{Frob}}
\providecommand{\GL}{}\renewcommand{\GL}{\operatorname{GL}}
\providecommand{\Gal}{}\renewcommand{\Gal}{\operatorname{Gal}}
\providecommand{\Hens}{}\renewcommand{\Hens}{\operatorname{Hens}}
\providecommand{\Hilb}{}\renewcommand{\Hilb}{\operatorname{Hilb}}
\providecommand{\Hom}{}\renewcommand{\Hom}{\operatorname{Hom}}
\providecommand{\IndCons}{}\renewcommand{\IndCons}{\operatorname{IndCons}}
\providecommand{\Kar}{}\renewcommand{\Kar}{\operatorname{Kar}}
\providecommand{\LT}{}\renewcommand{\LT}{\operatorname{LT}}
\providecommand{\Lie}{}\renewcommand{\Lie}{\operatorname{Lie}}
\providecommand{\Loc}{}\renewcommand{\Loc}{\operatorname{Loc}}
\providecommand{\Mod}{}\renewcommand{\Mod}{\mathsf{Mod}}
\providecommand{\N}{}\renewcommand{\N}{\mathbb{N}}
\providecommand{\Needspace}{}\renewcommand{\Needspace}{\par}
\providecommand{\O}{}\renewcommand{\O}{{\mathcal{O}}}
\providecommand{\Ob}{}\renewcommand{\Ob}{\mathrm{Ob}}
\providecommand{\Oe}{}\renewcommand{\Oe}{\mathcal O_E}
\providecommand{\PP}{}\renewcommand{\PP}{\mathbf{P}}
\providecommand{\Perf}{}\renewcommand{\Perf}{\operatorname{Perf}}
\providecommand{\Pic}{}\renewcommand{\Pic}{\mathrm{Pic}}
\providecommand{\Pm}{}\renewcommand{\Pm}{\mathbf P^{m}}
\providecommand{\Pone}{}\renewcommand{\Pone}{\mathbf P^{1}}
\providecommand{\Pqn}{}\renewcommand{\Pqn}{\mathbf P^{n}}
\providecommand{\ProFin}{}\renewcommand{\ProFin}{\operatorname{ProFin}}
\providecommand{\Proet}{}\renewcommand{\Proet}{\mathrm{pro\acute{e}t}}
\providecommand{\Q}{}\renewcommand{\Q}{\mathbb Q}
\providecommand{\Qb}{}\renewcommand{\Qb}{\overline{\mathbf Q}_{\ell}}
\providecommand{\Qbar}{}\renewcommand{\Qbar}{\overline{\mathbf Q}}
\providecommand{\QbarX}{}\renewcommand{\QbarX}{\overline{\mathbf Q}_{\ell,X}}
\providecommand{\Qbarl}{}\renewcommand{\Qbarl}{{\overline{\mathbf Q}_{\ell}}}
\providecommand{\Qell}{}\renewcommand{\Qell}{{\mathbf{Q}_{\ell}}}
\providecommand{\Ql}{}\renewcommand{\Ql}{\mathbf Q_{\ell}}
\providecommand{\Qlb}{}\renewcommand{\Qlb}{{\overline{\mathbf Q}_\ell}}
\providecommand{\Qlbar}{}\renewcommand{\Qlbar}{\overline{\mathbf Q}_\ell}
\providecommand{\R}{}\renewcommand{\R}{\mathbb{R}}
\providecommand{\RG}{}\renewcommand{\RG}{\operatorname{\mathrm{R}\Gamma}}
\providecommand{\RGamma}{}\renewcommand{\RGamma}{R\Gamma}
\providecommand{\RGammac}{}\renewcommand{\RGammac}{R\Gamma_c}
\providecommand{\RHom}{}\renewcommand{\RHom}{R\mathcal{H}om}
\providecommand{\RHomX}{}\renewcommand{\RHomX}{R\mathcal H om_X}
\providecommand{\Rat}{}\renewcommand{\Rat}{\operatorname{IsRat}}
\providecommand{\Res}{}\renewcommand{\Res}{\operatorname{Res}}
\providecommand{\Rlim}{}\renewcommand{\Rlim}{R\!\varprojlim}
\providecommand{\SIm}{}\renewcommand{\SIm}{\operatorname{SIm}}
\providecommand{\Sch}{}\renewcommand{\Sch}{\mathsf{Sch}}
\providecommand{\Sh}{}\renewcommand{\Sh}{\operatorname{Sh}}
\providecommand{\Span}{}\renewcommand{\Span}{\operatorname{span}}
\providecommand{\Spec}{}\renewcommand{\Spec}{\operatorname{Spec}}
\providecommand{\Tr}{}\renewcommand{\Tr}{\operatorname{Tr}}
\providecommand{\Trsm}{}\renewcommand{\Trsm}{\operatorname{Tr}^{\mathrm{sm}}}
\providecommand{\Wt}{}\renewcommand{\Wt}{\operatorname{Wt}}
\providecommand{\Z}{}\renewcommand{\Z}{\mathbb{Z}}
\providecommand{\Zar}{}\renewcommand{\Zar}{\mathrm{zar}}
\providecommand{\Zell}{}\renewcommand{\Zell}{{\mathbf{Z}_{\ell}}}
\providecommand{\Zl}{}\renewcommand{\Zl}{{\mathbf Z_\ell}}
\providecommand{\affet}{}\renewcommand{\affet}[1]{{#1}^{\mathrm{aff}}_{\mathrm{ét}}}
\providecommand{\affproet}{}\renewcommand{\affproet}[1]{{#1}^{\mathrm{aff}}_{\mathrm{proét}}}
\providecommand{\boxt}{}\renewcommand{\boxt}{\boxtimes_A^{\mathrm L}}
\providecommand{\cF}{}\renewcommand{\cF}{\mathcal{F}}
\providecommand{\cHom}{}\renewcommand{\cHom}{R\mathcal{H}\!om}
\providecommand{\cI}{}\renewcommand{\cI}{\mathcal{I}}
\providecommand{\cL}{}\renewcommand{\cL}{\mathcal L}
\providecommand{\cM}{}\renewcommand{\cM}{\mathcal M}
\providecommand{\cN}{}\renewcommand{\cN}{\mathcal N}
\providecommand{\cO}{}\renewcommand{\cO}{\mathcal O}
\providecommand{\cT}{}\renewcommand{\cT}{\mathcal T}
\providecommand{\checked}{}\renewcommand{\checked}{\color{red}\aftergroup\restorecolor}
\providecommand{\coeff}{}\renewcommand{\coeff}{\operatorname{coeff}}
\providecommand{\coker}{}\renewcommand{\coker}{\operatorname{coker}}
\providecommand{\colim}{}\renewcommand{\colim}{\operatorname{colim}}
\providecommand{\constructedby}{}\renewcommand{\constructedby}[1]{\par\smallskip\noindent\textcolor{purple!60!black}{\textsf{Constructed by:} }
  \footnotesize\sloppy\nolinkurl{#1}\par}
\providecommand{\deepresult}{}\renewcommand{\deepresult}[1]{\par\smallskip\noindent\textcolor{orange!45!black}{
  \textsf{Blueprint role: named deep theorem node; see the cluster proof architecture.}}
  \par\noindent\textcolor{orange!45!black}{\textsf{Mathematical source:} }#1\par}
\providecommand{\dt}{}\renewcommand{\dt}{\mathbin{\otimes^{\mathbf L}}}
\providecommand{\dw}{}\renewcommand{\dw}{\operatorname{dw}}
\providecommand{\et}{}\renewcommand{\et}[1]{{#1}_{\mathrm{ét}}}
\providecommand{\ev}{}\renewcommand{\ev}{\operatorname{ev}}
\providecommand{\fib}{}\renewcommand{\fib}{\operatorname{fib}}
\providecommand{\formalizationnote}{}\renewcommand{\formalizationnote}[1]{\begin{quote}\small\raggedright\textbf{Formalization note.} #1\end{quote}}
\providecommand{\fpqc}{}\renewcommand{\fpqc}{X_{\mathrm{fpqc}}}
\providecommand{\gammacont}{}\renewcommand{\gammacont}{\Gamma_{\mathrm{cont}}}
\providecommand{\gm}{}\renewcommand{\gm}{\mathbf G_m}
\providecommand{\hocolim}{}\renewcommand{\hocolim}{\operatorname{hocolim}}
\providecommand{\holim}{}\renewcommand{\holim}{\operatorname{holim}}
\providecommand{\hrB}{}\renewcommand{\hrB}{\@ifnextchar]{\hrEnd}{\hrStep}}
\providecommand{\hrEnd}{}\renewcommand{\hrEnd}{\@firstoftwo{}}
\providecommand{\hrStep}{}\renewcommand{\hrStep}{\@firstoftwo{\hrB}}
\providecommand{\id}{}\renewcommand{\id}{\operatorname{id}}
\providecommand{\kunFl}{}\renewcommand{\kunFl}{\mathbf F_\ell}
\providecommand{\kunZn}{}\renewcommand{\kunZn}{\mathbf Z/n\mathbf Z}
\providecommand{\kuncd}{}\renewcommand{\kuncd}{\operatorname{cd}}
\providecommand{\kundim}{}\renewcommand{\kundim}{\dim}
\providecommand{\kunotimes}{}\renewcommand{\kunotimes}[1]{\otimes^{\mathbf L}_{#1}}
\providecommand{\kunsupp}{}\renewcommand{\kunsupp}{\operatorname{Supp}}
\providecommand{\lotimes}{}\renewcommand{\lotimes}{\otimes^{\mathrm L}}
\providecommand{\m}{}\renewcommand{\m}{{\mathfrak{m}}}
\providecommand{\n}{}\renewcommand{\n}{{\mathfrak{n}}}
\providecommand{\nagataGap}{}\renewcommand{\nagataGap}[1]{\par\noindent\textbf{Formalization gap.} #1\par}
\providecommand{\openobligation}{}\renewcommand{\openobligation}[1]{\par\smallskip\noindent\textcolor{red!65!black}{\textsf{Open obligation:} }
  \footnotesize\sloppy\nolinkurl{#1}\par}
\providecommand{\p}{}\renewcommand{\p}{{\mathfrak{p}}}
\providecommand{\pr}{}\renewcommand{\pr}{\operatorname{pr}}
\providecommand{\preshv}{}\renewcommand{\preshv}[1]{\mathrm{PreShv}(#1)}
\providecommand{\pro}{}\renewcommand{\pro}[1]{\mathrm{Pro}(#1)}
\providecommand{\proet}{}\renewcommand{\proet}{X_{\mathrm{proét}}}
\providecommand{\q}{}\renewcommand{\q}{{\mathfrak{q}}}
\providecommand{\rad}{}\renewcommand{\rad}{\operatorname{rad}}
\providecommand{\rank}{}\renewcommand{\rank}{\operatorname{rank}}
\providecommand{\red}{}\renewcommand{\red}{\color{red}}
\providecommand{\restorecolor}{}\renewcommand{\restorecolor}{\color{black}}
\providecommand{\retiredobligation}{}\renewcommand{\retiredobligation}[1]{\par\smallskip\noindent\textcolor{red!75!black}{\textsf{Retired invalid route:} }
  \footnotesize\sloppy\nolinkurl{#1}\par}
\providecommand{\rk}{}\renewcommand{\rk}{\operatorname{rk}}
\providecommand{\sep}{}\renewcommand{\sep}{\textnormal{sep}}
\providecommand{\shv}{}\renewcommand{\shv}[1]{\mathrm{Shv}(#1)}
\providecommand{\spec}{}\renewcommand{\spec}[1]{\mathrm{Spec}(#1)}
\providecommand{\stacks}{}\renewcommand{\stacks}[1]{\href{https://stacks.math.columbia.edu/tag/#1}{\textnormal{Stacks, Tag~#1}}}
\providecommand{\stacksproject}{}\renewcommand{\stacksproject}[1]{{\cite[\href{https://stacks.math.columbia.edu/tag/#1}{Tag #1}]{stacks-project}}}
\providecommand{\statusnote}{}\renewcommand{\statusnote}[1]{\begin{quote}\small\raggedright\textbf{Status.} #1\end{quote}}
\providecommand{\tildeideal}{}\renewcommand{\tildeideal}[2]{\widetilde{(#1, #2)}}
\providecommand{\tildering}{}\renewcommand{\tildering}[3]{{#1}_{\widetilde{(#2, #3)}}}
\providecommand{\tr}{}\renewcommand{\tr}{\operatorname{tr}}
\providecommand{\wt}{}\renewcommand{\wt}{\operatorname{wt}}

\begin{document}
\pagestyle{plain}
\begin{center}
{\LARGE\bfseries Weil Conjectures\par}
\vspace{0.7em}
{\large Chapter 14: Purity of real sheaves\par}
\vspace{0.5em}
{\small\color{bpmuted}Snapshot 2026-09-26, revision 8e4cb88e2a\par}
\end{center}
\vspace{0.5em}
{\small\tableofcontents}
\bigskip\hrule\bigskip

% object: chap-Purity-of-real-sheaves (blueprint/chap-Purity-of-real-sheaves/_section.md)
\section*{\texorpdfstring{14\quad Purity of real sheaves}{14 Purity of real sheaves}}
\phantomsection\label{bp1}%
\addcontentsline{toc}{section}{\texorpdfstring{14\quad Purity of real sheaves}{14 Purity of real sheaves}}
% text: prose-14-01
\phantomsection\label{bp2}%
In Chapter \bpout{17} the Fourier transform of a pure sheaf produces a lisse sheaf \(M'\) whose Frobenius eigenvalues we cannot compute. What we do know is that a twist of \(M'\) is a quotient of a sheaf whose traces are real, and from this alone we must conclude that \(M'\) is pure. In Chapter \bpout{18} the same need returns for the direct images of a family of curves, which we can make real but do not know to be mixed. Both arguments rest on one theorem (Theorem \hyperref[bp82]{14.26}):

\begin{quote}
on a smooth geometrically connected curve, every irreducible constituent of a real lisse Weil sheaf is pure, of weight equal to its determinant weight.
\end{quote}

Before this theorem we know the weights of one kind of sheaf only: a lisse sheaf of rank one on a curve has finite geometric monodromy, and therefore it is pure (Section \hyperref[bp3]{14.1}). In particular the determinant of every lisse sheaf is pure, and its weight divided by the rank, the determinant weight, is the average weight of the eigenvalues at each point. The theorem says that on an irreducible constituent of a real sheaf no eigenvalue deviates from this average.

The tool is the trace formula. On an affine curve \(U_0\) it expresses the \(L\)-function of a lisse sheaf \(\mathcal G\) as

\[
L(U_0,\mathcal G,t)
=\prod_{x\in|U_0|}\det(1-\operatorname{Frob}_xt^{\deg(x)}\mid\mathcal G_{\bar x})^{-1}
=\frac{\det(1-Ft\mid H^1_c(U,\mathcal G))}{\det(1-Ft\mid H^2_c(U,\mathcal G))}.
\]

The poles of the right-hand side are controlled by \(H^2_c\), which is easy to compute (Lemma \hyperref[bp66]{14.20}). On the left, a large eigenvalue at a point \(x\) gives a pole of the local factor at \(x\) close to the origin. But this pole could be cancelled by a zero of the product of the other factors, as in \((1-2t)^{-1}\cdot(1-2t)=1\). Positive factors, however, cannot cancel in this way. Indeed, suppose all factors are power series with nonnegative coefficients and constant term \(1\). Then the coefficients of each factor are bounded by those of the product, and Lemma \hyperref[bp75]{14.24} shows that no factor can converge on a smaller disc than the product.

This idea is Rankin's. In 1939 he improved the bound for Ramanujan's function \(\tau(n)\) to \(O(n^{29/5})\) by studying the Dirichlet series \(\sum_n\tau(n)^2n^{-s}\), whose coefficients are nonnegative. Deligne used the same argument in his first proof [\hyperref[bp84]{Del74}].

For the Euler factors to be positive, the traces must be real. We therefore call a lisse Weil sheaf real if its local characteristic polynomials have real coefficients under the fixed embedding \(\iota\) of Convention \bpslug{situation-{}\allowbreak{}weil} (Definition \hyperref[bp69]{14.21}). For example, every sheaf with rational characteristic polynomials is real, such as a Tate twist of the constant sheaf. On an even tensor power \(\mathcal G^{\otimes2m}\) of a real sheaf the traces are \(2m\)th powers of real numbers, and this makes the Euler factors positive. From this, Proposition \hyperref[bp78]{14.25} bounds each local eigenvalue \(\alpha\) of \(\mathcal G\) by \(|\iota(\alpha)|^2\le q_x^\beta\), where \(\beta\) is the largest determinant weight of the constituents.

This is not yet the theorem, because a real sheaf need not be pure. For example, the sheaf \(\overline{\mathbf Q}_\ell\oplus\overline{\mathbf Q}_\ell(-1)\) is real, with constituents of weights \(0\) and \(2\), but the bound only gives \(2\) for both. To separate the constituents we pass to exterior powers. In \(\bigwedge^2(\overline{\mathbf Q}_\ell\oplus\overline{\mathbf Q}_\ell(-1))=\det\overline{\mathbf Q}_\ell(-1)\otimes\overline{\mathbf Q}_\ell\) the factor \(\overline{\mathbf Q}_\ell\) appears multiplied by the known eigenvalue \(q_x\) of \(\det\overline{\mathbf Q}_\ell(-1)\), so the bound \(2\) for this summand leaves the bound \(0\) for \(\overline{\mathbf Q}_\ell\) itself. In general, Proposition proposition-determinant-weights-tensor-operations finds the exterior power in which a given group of constituents carries the largest determinant weight, and the same argument applies there. Positivity, however, only ever gives upper bounds. The lower bound comes from applying the same argument to the dual, and in Chapter \bpout{19} this role is played by Poincaré duality.

% object: sec-Sheaves-of-rank-one (blueprint/chap-Purity-of-real-sheaves/sec-Sheaves-of-rank-one/_section.md), content hidden
\subsection*{\texorpdfstring{14.1\quad Sheaves of rank one}{14.1 Sheaves of rank one}}
\phantomsection\label{bp3}%
\addcontentsline{toc}{subsection}{\texorpdfstring{14.1\quad Sheaves of rank one}{14.1 Sheaves of rank one}}
% text: prose-14-02
\phantomsection\label{bp4}%
Rankin's argument bounds the local eigenvalues of a sheaf by its determinant weights, so it needs the determinant of every lisse sheaf to be pure (Corollary \hyperref[bp27]{14.8}). Since a determinant has rank one, we prove more generally that every rank-one lisse Weil sheaf on a smooth geometrically connected curve is pure (Corollary \hyperref[bp23]{14.6}), and the later sections Sections \hyperref[bp31]{14.2}–\hyperref[bp80]{Section 14.4} are built on this fact.

In rank one the question is about characters. Let \(U_0\) be a smooth geometrically connected curve and \(s\in W(U_0)\) a lift of Frobenius of degree one (Convention \bpslug{situation-{}\allowbreak{}weil}). Then a rank-one lisse Weil sheaf is a character \(\chi:W(U_0)\to E^\times\), continuous on \(\pi_1(U)\). If \(\chi\) is trivial on \(\pi_1(U)\), the sheaf is one of the sheaves \(\mathcal B_b\) with \(b=\chi(s)\); at a closed point \(x\) its eigenvalue is \(b^{\deg(x)}\), so it is pure of weight \(2\log_q|\iota(b)|\). A Kummer sheaf is the next simplest case. If \(n\mid q-1\), the cover \(t\mapsto t^n\) of \(\mathbf G_m\) is a \(\mu_n\)-torsor, and so a character of \(\mu_n(k_0)\) gives a rank-one sheaf on \(\mathbf G_{m,k_0}\) whose geometric monodromy is cyclic of order dividing \(n\). Its eigenvalues are values of the character on \(\mu_n(k_0)\), hence roots of unity, and it is pure of weight \(0\).

Up to a twist, every rank-one sheaf behaves like these two. Indeed, by Theorem \hyperref[bp20]{14.5} the image \(\chi(\pi_1(U))\) is finite. It follows that \(\chi(\operatorname{Frob}_x)\) is \(\chi(s)^{\deg(x)}\) times a root of unity, so the sheaf is pure of weight \(2\log_q|\iota(\chi(s))|\).

Finiteness is a statement about Weil sheaves, because on the geometric curve there are rank-one local systems with infinite monodromy. For example, the tame fundamental group of \(\mathbf G_{m,k}\) has the quotient \(\mathbf Z_\ell(1)\simeq\mathbf Z_\ell\), and \(a\mapsto(1+\ell)^a\) is a character of infinite order. But it has no Weil structure. Indeed, a Weil structure forces \(\chi(s^{-1}gs)=\chi(g)\) for \(g\in\pi_1(U)\), while conjugation by \(s\) acts on \(\mathbf Z_\ell(1)\) by multiplication by a power \(q^{\pm1}\). A character of \(\mathbf Z_\ell(1)\) that extends to the Weil group therefore satisfies \(\chi(g)^{q-1}=1\), and it has finite order.

The general proof follows the same pattern. Since there is no explicit quotient to compute with, it uses a logarithm instead: a power \(\chi^N\) takes values in \(1+\mathfrak m_E^r\), where the \(\ell\)-adic logarithm converges, and then \(\log\chi^N\) is an additive character of \(\pi_1(U)\) invariant under conjugation by \(s\), that is, a class in \(H^1(U,E)\) fixed by Frobenius (Lemma \hyperref[bp17]{14.4}). Such a class is zero (Proposition \hyperref[bp13]{14.3}), so \(\chi^N\) is trivial on \(\pi_1(U)\).

For a proper curve, Frobenius fixes no class in \(H^1\) because the Jacobian has only finitely many rational points. If \(C_0\) is a proper curve with Jacobian \(J_0\), Kummer theory and Poincaré duality identify \(\varprojlim_nH^1(C,\mathbf Z/\ell^n)\) with the dual of the Tate module \(T_\ell J\) (Proposition \hyperref[bp10]{14.2}). A fixed vector in the dual therefore means that \(1\) is an eigenvalue of Frobenius on \(V_\ell J=T_\ell J\otimes\mathbf Q_\ell\), and after scaling, Frobenius would then fix a nonzero element of \(T_\ell J\). Such an element is a compatible system of \(\ell\)-power torsion points of \(J\) defined over \(k_0\), of unbounded order. These points cannot all lie in the finite group \(J_0(k_0)\), and so there is no such element (Proposition \bpout{B.15}).

For an open curve, the residue sequence contains in addition the boundary term \(H^0(D,E(-1))\). A power of Frobenius acts on it by a power of \(q\ne1\), so it has no fixed vector either. For example, for \(\mathbf G_m\) this is the computation above, since there \(H^1(\mathbf G_{m,k},\mathbf Q_\ell)=\mathbf Q_\ell(-1)\), with Frobenius acting by \(q\).

% object: subsec-Frobenius-invariants (blueprint/chap-Purity-of-real-sheaves/sec-Sheaves-of-rank-one/subsec-Frobenius-invariants/_section.md), content hidden
\subsubsection*{\texorpdfstring{14.1.1\quad Frobenius invariants}{14.1.1 Frobenius invariants}}
\phantomsection\label{bp5}%
\addcontentsline{toc}{subsubsection}{\texorpdfstring{14.1.1\quad Frobenius invariants}{14.1.1 Frobenius invariants}}
% text: prose-14-03
\phantomsection\label{bp6}%
Over \(\mathbf C\) the link between \(H^1\) of a curve and its Jacobian is classical. There the Jacobian is a torus \(V/\Lambda\) with \(\Lambda=H_1(C,\mathbf Z)\), and \(H^1(C,\mathbf Z)=\Hom(\Lambda,\mathbf Z)\). Proposition \hyperref[bp10]{14.2} is the \(\ell\)-adic version, with \(T_\ell J\) in place of \(\Lambda\otimes\mathbf Z_\ell\). Here the link with the Jacobian is Kummer theory, which is a statement about \(\mu_{\ell^n}\) and line bundles. For this reason we work with the coefficients \(\mathbf Z/\ell^n\) and pass to \(E\) only at the end. Apart from this, the only arithmetic input is the finiteness of \(J_0(k_0)\), which holds because \(J_0\) is of finite type over a finite field.

% object: lemma-proper-smooth-model (blueprint/chap-Purity-of-real-sheaves/sec-Sheaves-of-rank-one/subsec-Frobenius-invariants/lemma-proper-smooth-model/_section.md)
\phantomsection\label{bp7}%
\begin{bpstatementit}{Lemma~14.1}{ (A proper smooth model)}
A smooth geometrically connected curve \(U_0/k_0\) has a smooth projective geometrically connected compactification \(j_0:U_0\hookrightarrow C_0\). The reduced boundary \(D_0\) is finite étale over \(k_0\).
\end{bpstatementit}

% proof: proof/lemma-proper-smooth-model
\begin{proof}[{Proof}]\phantomsection\label{bp8}%
We compactify by Nagata's theorem and then normalize. On a curve, normality is the same as regularity, and over a perfect field regularity is smoothness, so the normalization is the model we want.

By Theorem \bpout{A.2}, \(U_0\) is an open subscheme of a proper \(k_0\)-scheme, and we take the reduced closure of \(U_0\) in it and normalize this closure in \(k_0(U_0)\).  Since a scheme of finite type over a field is Nagata, the normalization is finite [Stacks, \href{https://stacks.math.columbia.edu/tag/035B}{035B}, \href{https://stacks.math.columbia.edu/tag/035S}{035S}], and it is an isomorphism over the normal open \(U_0\).  The resulting proper normal curve is regular, because a one-dimensional noetherian normal local domain is a discrete valuation ring. Hence it is smooth over the perfect field \(k_0\) [Stacks, \href{https://stacks.math.columbia.edu/tag/0B8Y}{0B8Y}].  Its boundary is zero-dimensional and proper over \(k_0\), hence finite. Since \(k_0\) is perfect, the reduced boundary is étale.  After base change to \(k\), the curve \(U\) is dense and connected, and therefore its closure \(C\) is connected.  Finally, a proper scheme of dimension one over a field is projective [Stacks, \href{https://stacks.math.columbia.edu/tag/0A26}{0A26}].
\end{proof}

% text: prose-14-04
\phantomsection\label{bp9}%
Once we have a smooth projective model \(C_0\), the Jacobian \(J_0\) of \(C_0\) exists (Appendix \bpout{C}), and we can compare the cohomology of \(C\) with its Tate module. We first make the comparison with finite coefficients, where Kummer theory applies directly. The proposition ends with a matrix description of Frobenius, which is the form used in Proposition \hyperref[bp13]{14.3}.

% object: proposition-kummer-duality-integral-tate-module (blueprint/chap-Purity-of-real-sheaves/sec-Sheaves-of-rank-one/subsec-Frobenius-invariants/proposition-kummer-duality-integral-tate-module/_section.md)
\phantomsection\label{bp10}%
\begin{bpstatementit}{Proposition~14.2}{ (Kummer, duality and the integral Tate module)}
Let \(C_0/k_0\) be a smooth proper geometrically connected curve with Jacobian \(J_0\), and put \(A_n=\mathbf Z/\ell^n\).  There are canonical compatible Galois-equivariant identifications

\[
H^1(C,A_n(1))\simeq J[\ell^n](k),\qquad
\varprojlim_n H^1(C,A_n)
\simeq\Hom_{\mathbf Z_\ell}(T_\ell J,\mathbf Z_\ell).
\]

Here the transition maps are multiplication by \(\ell\) on \(J[\ell^n](k)\) and reduction of coefficients on \(H^1(C,A_n)\).  If arithmetic Frobenius on \(T_\ell J\) has matrix \(B\), then geometric Frobenius on \(\Hom_{\mathbf Z_\ell}(T_\ell J,\mathbf Z_\ell)\) has matrix \(B^t\) in the dual basis.
\end{bpstatementit}

% proof: proof/proposition-kummer-duality-integral-tate-module
\begin{proof}[{Proof}]\phantomsection\label{bp11}%
By the Kummer sequence, \(H^1(C,\mu_{\ell^n})\) is the \(\ell^n\)-torsion of the Picard group, and Poincaré duality turns this into a statement about the untwisted \(H^1\).

The first identification is Proposition \bpout{C.17}, which also gives its Galois equivariance, and it shows that the coefficient map \(z\mapsto z^\ell\) from \(\mu_{\ell^{n+1}}\) to \(\mu_{\ell^n}\) corresponds to \([\ell]\) on the Jacobian.

For the dual Tate module, Poincaré duality at finite level pairs the untwisted \(H^1\) with the twisted one that we have just identified with \(J[\ell^n](k)\). Derived duality gives a statement about the single group \(H^1\) because all the groups involved are free \(A_n\)-modules. Indeed, we have \(H^0(C,A_n)=A_n\) and \(H^2(C,A_n(1))=A_n\) by Proposition \bpout{D.2}, while the group \(H^1(C,A_n(1))\) is free of rank \(2g\) by the identification above and Corollary \bpout{C.15}. Since \(A_n(1)\simeq A_n\) over \(k\), the untwisted \(H^1\) is free of rank \(2g\) as well. Finally, a bounded complex with projective cohomology modules is the sum of its shifted cohomology modules in the derived category, so \(A_n\)-duality acts degree by degree. In this situation derived Poincaré duality, Theorem \bpout{10.12}, applied to the constant sheaf, induces in degree one an isomorphism

\[
H^1(C,A_n)\xrightarrow{\sim}
\Hom_{A_n}(H^1(C,A_n(1)),A_n).
\]

This map is the adjoint of cup product followed by the trace map of the duality theorem, and both cup product and trace commute with reduction of coefficients and with the action of \(\operatorname{Gal}(k/k_0)\).  Because of the twist on the second factor, Frobenius acts trivially on the target \(A_n\).

In the limit, the finite-level dualities combine into a duality with the Tate module. Since the transition maps \([\ell]\) are surjective, we have \(J[\ell^n](k)=T_\ell J/\ell^nT_\ell J\). Moreover \(T_\ell J\) is finite free by Proposition \bpout{B.14}, and a homomorphism of finite free \(\mathbf Z_\ell\)-modules is determined by its reductions modulo all \(\ell^n\), and every compatible system of reductions lifts. Hence the inverse limit of the adjoints at finite level is the dual of \(T_\ell J\). We obtain a perfect pairing

\[
\varprojlim_nH^1(C,A_n)\times T_\ell J\longrightarrow\mathbf Z_\ell,
\]

invariant under Frobenius.  Since geometric Frobenius acts on \(T_\ell J\) by \(B^{-1}\), its matrix on \(\varprojlim_nH^1(C,A_n)\) in the dual basis is \((B^{-1})^{-t}=B^t\).
\end{proof}

% text: prose-14-05
\phantomsection\label{bp12}%
For the elliptic curve \(y^2=x^3-x\) over \(\mathbf F_3\) of the introduction, arithmetic Frobenius \(B\) has eigenvalues \(\pm i\sqrt3\) on \(T_\ell J\). By the proposition, geometric Frobenius acts on \(\varprojlim_nH^1(C,A_n)\) with the same eigenvalues. Since \(\det(B^t-1)=4\ne0\), it has no fixed vector. The next proposition proves the same for every smooth curve, including open ones.

% object: proposition-no-fixed-class-degree-one (blueprint/chap-Purity-of-real-sheaves/sec-Sheaves-of-rank-one/subsec-Frobenius-invariants/proposition-no-fixed-class-degree-one/_section.md)
\phantomsection\label{bp13}%
\begin{bpstatementit}{Proposition~14.3}{ (No fixed class in degree one)}
For a smooth geometrically connected curve \(U_0/k_0\), geometric Frobenius has no fixed vector on \(H^1(U,E)\) for every finite \(E/\mathbf Q_\ell\).  More precisely,

\[
\left(\varprojlim_nH^1(U,\mathbf Z/\ell^n)\right)^{F=1}=0.
\]
\end{bpstatementit}

% proof: proof/proposition-no-fixed-class-degree-one
\begin{proof}[{Proof}]\phantomsection\label{bp14}%
On the compactification the vanishing holds because \(J_0(k_0)\) is finite, and on the boundary it holds because Frobenius acts by a power of \(q\), which is not \(1\). The residue sequence reduces the open case to the proper one, and at the end we pass from integral to rational coefficients.

\emph{The proper case.} Let \(U_0=C_0\) be proper, with Jacobian \(J_0\). Here the cohomology is the dual of the Tate module, and the question is whether \(1\) is an eigenvalue of Frobenius on it.  By Proposition \bpout{B.15}, arithmetic Frobenius has no fixed vector on \(T_\ell J\); indeed, a fixed compatible tuple would consist of \(\ell\)-power torsion points of the finite group \(J_0(k_0)\). This means that \(B-1\) is invertible on \(V_\ell J\).  On the other side, geometric Frobenius acts on the untwisted \(H^1\) by \(B^t\) (Proposition \hyperref[bp10]{14.2}). Since \(\det(B^t-1)=\det(B-1)\ne0\), it fixes neither a nonzero rational vector nor a nonzero integral compatible tuple.  The argument applies to every power of Frobenius, so the same holds over every finite extension of \(k_0\), which we need in the open case.

\emph{The open case.} A fixed class has zero residues at the boundary points, so it extends to the compactification, where the proper case shows that it is zero. To make this precise, choose \(j:U\hookrightarrow C\) as in Lemma \hyperref[bp7]{14.1}, with boundary \(i:D\hookrightarrow C\). At each finite level there is then a canonical residue sequence

\[
0\longrightarrow H^1(C,A_n)\longrightarrow H^1(U,A_n)
\longrightarrow H^0(D,A_n(-1)).
\]

To obtain it with the correct twist and Frobenius action, we use the dualizing complexes.  By transitivity, \(i^!\omega_C=\omega_D=A_n\). Since \(\omega_C=A_n(1)[2]\), it follows that \(i^!A_{n,C}=A_{n,D}(-1)[-2]\).  The localization triangle \(i_*i^!A_{n,C}\to A_{n,C}\to Rj_*A_{n,U}\) now yields the displayed sequence, and all its maps are natural in the coefficients and under automorphisms of the model.

Choose \(m>0\) such that \(F^m\) fixes all geometric boundary points, and put \(Q=q^m\).  On \(\varprojlim_nH^0(D,A_n(-1))\), which is finite free over \(\mathbf Z_\ell\), the operator \(F^m\) acts as \(Q\). Since \(Q-1\ne0\), it has no nonzero fixed element. A fixed compatible class on \(U\) therefore has zero residues at every level, so by exactness it lifts uniquely to \(H^1(C,A_n)\). By uniqueness the lifts are compatible and fixed by \(F^m\), and the proper case over \(\mathbf F_Q\) shows that they vanish.

\emph{Rational coefficients.} To pass from compatible systems of classes to \(H^1(U,E)\), we identify \(H^1(U,\mathbf Q_\ell)\) with the rationalization of the inverse limit and check that rationalizing does not create fixed vectors. By Proposition \bpout{2.55}, applied to the constant sheaves \(A_n\), we have \(R\Gamma(U,\mathbf Z_\ell)=R\varprojlim_nR\Gamma(U,A_n)\), and hence an exact sequence

\[
0\to\varprojlim\nolimits^1_nH^0(U,A_n)\to H^1(U,\mathbf Z_\ell)
\to\varprojlim_nH^1(U,A_n)\to0.
\]

The \(\varprojlim^1\)-term vanishes because the groups \(H^0(U,A_n)=A_n\) have surjective transition maps.  On \(U_{\mathrm{pro}\acute{e}t}\) the sheaf \(\mathbf Q_\ell\) is \(\mathbf Z_\ell[1/\ell]\) (Lemma \bpout{2.51}), a filtered colimit of copies of \(\mathbf Z_\ell\), and cohomology on the quasi-compact and quasi-separated scheme \(U\) commutes with filtered colimits of sheaves (Corollary \bpout{2.29}, applied to sheaves placed in degree zero).  So \(H^1(U,\mathbf Q_\ell)=H^1(U,\mathbf Z_\ell)\otimes\mathbf Q_\ell\) is the rationalization of the inverse limit.

The inverse limit of the residue sequences is exact on the left, and rationalization is exact, so we obtain an exact sequence of finite-dimensional \(\mathbf Q_\ell\)-spaces \(0\to H^1(C,\mathbf Q_\ell)\to H^1(U,\mathbf Q_\ell)\to I\to0\), where \(I\) is the image of the residue map inside the rationalized boundary term. By the proper case, \(F^m-1\) is invertible on \(H^1(C,\mathbf Q_\ell)\). On the boundary term, \(F^m-1\) acts by \(Q-1\), so it is injective on the \(F^m\)-stable subspace \(I\). Now let \(v\) be a vector of \(H^1(U,\mathbf Q_\ell)\) killed by \(F^m-1\). Its image in \(I\) is killed by \(F^m-1\), hence zero. Thus \(v\) comes from \(H^1(C,\mathbf Q_\ell)\), where it is zero by the proper case. Since a vector fixed by \(F\) is fixed by \(F^m\), it follows that Frobenius has no fixed rational vector, and extending scalars to \(E\) gives the assertion over \(E\).
\end{proof}


% object: subsec-Finite-monodromy (blueprint/chap-Purity-of-real-sheaves/sec-Sheaves-of-rank-one/subsec-Finite-monodromy/_section.md), content hidden
\subsubsection*{\texorpdfstring{14.1.2\quad Finite monodromy}{14.1.2 Finite monodromy}}
\phantomsection\label{bp15}%
\addcontentsline{toc}{subsubsection}{\texorpdfstring{14.1.2\quad Finite monodromy}{14.1.2 Finite monodromy}}
% text: prose-14-06
\phantomsection\label{bp16}%
With Proposition \hyperref[bp13]{14.3} in hand, the logarithm argument sketched at the start of Section \hyperref[bp3]{14.1} proves Theorem \hyperref[bp20]{14.5}. Corollaries \hyperref[bp23]{Corollary 14.6} and \hyperref[bp27]{Corollary 14.8} then translate finiteness of monodromy into weights. For a lisse sheaf of higher rank, however, the determinant weight (Definition \hyperref[bp26]{14.7}) is the only weight we can attach until Rankin's argument. Throughout, \(U_0\) is a smooth geometrically connected curve over \(k_0\) and \(s\in W(U_0)\) is a lift of degree one, as in Convention \bpslug{situation-{}\allowbreak{}weil}. In this setting, the first step identifies the additive characters to which the logarithm leads.

% object: lemma-invariant-additive-characters (blueprint/chap-Purity-of-real-sheaves/sec-Sheaves-of-rank-one/subsec-Finite-monodromy/lemma-invariant-additive-characters/_section.md)
\phantomsection\label{bp17}%
\begin{bpstatementit}{Lemma~14.4}{ (Invariant additive characters)}
Every continuous additive homomorphism \(a:\pi_1(U)\to E\) satisfying \(a(s^{-1}gs)=a(g)\) is zero.
\end{bpstatementit}

% proof: proof/lemma-invariant-additive-characters
\begin{proof}[{Proof}]\phantomsection\label{bp18}%
After scaling, the coordinates of \(a\) reduce modulo \(\ell^n\) to a compatible system of Frobenius-fixed classes in \(H^1(U,\mathbf Z/\ell^n)\), which vanish by Proposition \hyperref[bp13]{14.3}. This needs \(H^1\) to be identified with homomorphisms.

Recall that for a finite abelian group \(G\) the group \(H^1(U,G)\) classifies \(G\)-torsors. By [Stacks, \href{https://stacks.math.columbia.edu/tag/0DV5}{0DV5}], finite locally constant sheaves on the connected scheme \(U\) correspond to finite continuous \(\pi_1(U)\)-sets, and under this correspondence a \(G\)-torsor becomes a finite \(\pi_1(U)\)-set with a commuting simply transitive action of \(G\), that is, a continuous homomorphism \(\pi_1(U)\to G\). Hence \(H^1(U,G)=\Hom_{\mathrm{cont}}(\pi_1(U),G)\). In this description Frobenius acts on \(H^1(U,G)\) through the action of \(s\) on \(\pi_1(U)\) by conjugation. The Frobenius-fixed classes are therefore the homomorphisms \(b\) with \(b(s^{-1}gs)=b(g)\).

Choose a \(\mathbf Q_\ell\)-basis of \(E\), with coordinate maps \(t_i\).  Since \(\pi_1(U)\) is compact, each \(t_i a(\pi_1(U))\) is bounded, and after multiplying all coordinates by one power of \(\ell\) we may assume that they take values in \(\mathbf Z_\ell\).  Each coordinate is then a continuous homomorphism \(\pi_1(U)\to\mathbf Z_\ell\) invariant under conjugation by \(s\). Its reductions modulo \(\ell^n\) are compatible classes in \(H^1(U,\mathbf Z/\ell^n)\) fixed by Frobenius. By the proposition they vanish, and hence every coordinate of \(a\) is zero.
\end{proof}

% text: prose-14-07
\phantomsection\label{bp19}%
The invariance hypothesis of the lemma is automatic for characters of the Weil group. If \(\chi:W(U_0)\to E^\times\) is a character, then \(\chi(s^{-1}gs)=\chi(s)^{-1}\chi(g)\chi(s)=\chi(g)\), because \(E^\times\) is commutative. So the lemma applies as soon as \(\chi\) has been made additive on \(\pi_1(U)\).

% object: theorem-rank-one-sheaves-have-finite (blueprint/chap-Purity-of-real-sheaves/sec-Sheaves-of-rank-one/subsec-Finite-monodromy/theorem-rank-one-sheaves-have-finite/_section.md)
\phantomsection\label{bp20}%
\begin{bpstatementit}{Theorem~14.5}{ (Rank-one sheaves have finite geometric monodromy)}
Let \(U_0/k_0\) be a smooth geometrically connected curve, and let \(\chi:W(U_0)\to E^\times\) be the character of a rank-one lisse Weil sheaf.  Then \(\chi(\pi_1(U))\) is finite.
\end{bpstatementit}

% proof: proof/theorem-rank-one-sheaves-have-finite
\begin{proof}[{Proof}]\phantomsection\label{bp21}%
We show that a fixed power of \(\chi\) is trivial on \(\pi_1(U)\), by applying Lemma \hyperref[bp17]{14.4} to its logarithm. For this we need a subgroup \(1+\mathfrak m_E^r\) of finite index in \(\mathcal O_E^\times\) on which the logarithm is an injective homomorphism, and a power of \(\chi\) that lands in it.

\emph{Step 1: the logarithm.} The compact group \(\chi(\pi_1(U))\) lies in \(\mathcal O_E^\times\), because its image under the normalized valuation is a finite subgroup of \(\mathbf Z\) and therefore zero. Now write \(e=v_E(\ell)\) and choose an integer \(r>e/(\ell-1)\). We claim that \(\log\) is injective on \(1+\mathfrak m_E^r\). For \(x\in\mathfrak m_E^r\) and \(n\ge1\) we have \(v_E(n)=e\,v_\ell(n)\) and, by Legendre's formula, \(v_\ell(n!)\le(n-1)/(\ell-1)\), so that

\[
v_E(x^n/n)\ge nr-e\,v_\ell(n),
\qquad
v_E(x^n/n!)\ge r+(n-1)\bigl(r-e/(\ell-1)\bigr).
\]

Both bounds tend to infinity with \(n\) and are at least \(r\): for the second this is clear, and for the first it follows from \(n-1\ge\ell^{v_\ell(n)}-1\ge(\ell-1)v_\ell(n)\) and \((\ell-1)r>e\), which give \((n-1)r\ge e\,v_\ell(n)\). The logarithm and exponential series therefore converge on \(\mathfrak m_E^r\); the logarithm maps \(1+\mathfrak m_E^r\) into \(\mathfrak m_E^r\), and the exponential maps \(\mathfrak m_E^r\) into \(1+\mathfrak m_E^r\). By the usual identities of formal power series, \(\log:1+\mathfrak m_E^r\to\mathfrak m_E^r\) and \(\exp\) are mutually inverse continuous group homomorphisms.

\emph{Step 2: a power of \(\chi\).} The quotient \(\mathcal O_E^\times/(1+\mathfrak m_E^r)\) is finite, being isomorphic to \((\mathcal O_E/\mathfrak m_E^r)^\times\); let \(N>0\) be an integer that kills it. As \(\chi(\pi_1(U))\subset\mathcal O_E^\times\), we have \(\chi(g)^N\in1+\mathfrak m_E^r\) for \(g\in\pi_1(U)\). Hence \(a(g)=\log(\chi(g)^N)\) defines a continuous homomorphism \(\pi_1(U)\to E\), which is invariant under conjugation by \(s\) because \(\chi\) takes values in an abelian group.

\emph{Step 3: conclusion.} Lemma \hyperref[bp17]{14.4} gives \(a=0\), so \(\chi(g)^N=1\) for every \(g\in\pi_1(U)\), the logarithm being injective. The theorem follows, since \(T^N-1\) has only finitely many roots in \(E\).
\end{proof}

% text: prose-14-08
\phantomsection\label{bp22}%
Finiteness gives the sheaf a weight, and since the lift \(s\) is a choice, part of the proof checks that the weight does not depend on it.

% object: corollary-rank-one-weights (blueprint/chap-Purity-of-real-sheaves/sec-Sheaves-of-rank-one/subsec-Finite-monodromy/corollary-rank-one-weights/_section.md)
\phantomsection\label{bp23}%
\begin{bpstatementit}{Corollary~14.6}{ (Rank-one weights)}
Every rank-one lisse Weil sheaf on a smooth geometrically connected curve is punctually \(\iota\)-pure.  For \(c=\chi(s)\) its weight is \(\lambda_\iota(c)=2\log_q|\iota(c)|\).
\end{bpstatementit}

% proof: proof/corollary-rank-one-weights
\begin{proof}[{Proof}]\phantomsection\label{bp24}%
We separate from \(\chi\) the part that depends only on the degree. Put \(\varepsilon(w)=\chi(w)c^{-\deg w}\), a character of \(W(U_0)\) that is trivial on \(s\).  Every element of \(W(U_0)\) can be written as \(w=gs^d\) with \(g\in\pi_1(U)\) and \(d=\deg w\), and then \(\varepsilon(w)=\chi(g)\chi(s)^dc^{-d}=\chi(g)\). So \(\varepsilon\) has the same image as \(\chi|_{\pi_1(U)}\), which is finite by Theorem \hyperref[bp20]{14.5}.  At a closed point \(x\) the eigenvalue is \(\varepsilon(\operatorname{Frob}_x)c^{\deg(x)}\), where the first factor has finite order, hence complex absolute value one. Therefore the eigenvalue has absolute value \(q_x^{\lambda_\iota(c)/2}\).  This number does not depend on \(s\), because changing \(s\) multiplies \(c\) by a value of \(\chi\) on \(\pi_1(U)\), which has finite order.
\end{proof}

% text: prose-14-09
\phantomsection\label{bp25}%
For a Kummer sheaf on \(\mathbf G_{m,k_0}\) the corollary gives weight \(0\), since \(\chi(s)\) is then a root of unity. For \(\mathcal B_b\) it gives \(\lambda_\iota(b)\), the weight found in Proposition \bpout{13.7}.

Before purity is known, the eigenvalues of a lisse sheaf \(\mathcal V\) of rank \(r\) at a closed point are not controlled individually. Their product, however, is controlled, since it is the eigenvalue of the rank-one sheaf \(\det\mathcal V\), whose weight is given by Corollary \hyperref[bp23]{14.6}. Dividing by \(r\) we get the average weight of the eigenvalues, which is the only candidate for a weight of \(\mathcal V\). Indeed, if \(\mathcal V\) is pure of weight \(w\), then \(\det\mathcal V\) is pure of weight \(rw\), and the number defined below is \(w\). For a sheaf that is not pure, on the other hand, the number is only an average. For example, on \(\overline{\mathbf Q}_\ell\oplus\overline{\mathbf Q}_\ell(-1)\) the determinant is \(\overline{\mathbf Q}_\ell(-1)\), of weight \(2\), and the number is \(1\), the weight of no eigenvalue. For this reason we apply it to irreducible constituents: Proposition \hyperref[bp78]{14.25} bounds the eigenvalues of a real sheaf by the largest determinant weight of its constituents, and Theorem \hyperref[bp82]{14.26} then shows that each irreducible constituent is pure of weight equal to its determinant weight.

% object: definition-determinant-weight (blueprint/chap-Purity-of-real-sheaves/sec-Sheaves-of-rank-one/subsec-Finite-monodromy/definition-determinant-weight/_section.md)
\phantomsection\label{bp26}%
\begin{bpstatement}{Definition~14.7}{ (Determinant weight)}
Let \(\mathcal V\) be a nonzero lisse Weil sheaf of rank \(r\) on a smooth geometrically connected curve \(U_0\), let \(\rho\) be the corresponding representation of \(W(U_0)\), and let \(s\in W(U_0)\) be a lift of degree one. The \emph{determinant weight} of \(\mathcal V\) is

\[
\operatorname{dw}_\iota(\mathcal V)
=\frac{\lambda_\iota(\det\rho(s))}{r}
\]

(this does not depend on \(s\), Corollary \hyperref[bp27]{14.8}). For a lisse Weil sheaf \(\mathcal V\) we write \(\operatorname{DW}(\mathcal V)\) for the set of determinant weights of its irreducible constituents.
\end{bpstatement}

% object: corollary-determinant-weights (blueprint/chap-Purity-of-real-sheaves/sec-Sheaves-of-rank-one/subsec-Finite-monodromy/corollary-determinant-weights/_section.md)
\phantomsection\label{bp27}%
\begin{bpstatementit}{Corollary~14.8}{ (Determinant weights)}
Let \(\mathcal V\) be a nonzero lisse Weil sheaf of rank \(r\) on a smooth geometrically connected curve.  Its determinant has finite geometric monodromy and is punctually pure of weight \(r\operatorname{dw}_\iota(\mathcal V)\). The number \(\operatorname{dw}_\iota(\mathcal V)\) does not depend on the choice of \(s\) or of a finite coefficient field over which \(\mathcal V\) is defined.
\end{bpstatementit}

% proof: proof/corollary-determinant-weights
\begin{proof}[{Proof}]\phantomsection\label{bp28}%
The sheaf \(\det\mathcal V\) is a rank-one lisse Weil sheaf defined over a finite extension of \(E\), so Theorem \hyperref[bp20]{14.5} and Corollary \hyperref[bp23]{14.6} apply to it.  By the formula of Corollary \hyperref[bp23]{14.6}, with \(c=\det\rho(s)\), the weight of \(\det\mathcal V\) is \(r\operatorname{dw}_\iota(\mathcal V)\), and it does not depend on \(s\). Finally, a finite extension of scalars changes neither the character nor its absolute value under the fixed embedding.
\end{proof}

% text: prose-14-10
\phantomsection\label{bp29}%
The corollary applies to every irreducible constituent of \(\mathcal V\) and of every tensor or exterior power of \(\mathcal V\). Indeed, a finite-dimensional representation has a composition series, and each factor is again lisse and defined over a finite extension of \(\mathbf Q_\ell\), because after enlarging \(E\) the invariant subspaces and quotients have bases defined over \(E\).

% object: remark-determinant-weight-pure-constituents (blueprint/chap-Purity-of-real-sheaves/sec-Sheaves-of-rank-one/subsec-Finite-monodromy/remark-determinant-weight-pure-constituents/_section.md)
\phantomsection\label{bp30}%
\begin{bpstatement}{Remark~14.9}{ (Determinant weight of a sheaf with pure constituents)}
Let \(\mathcal V\) be a nonzero lisse Weil sheaf of rank \(r\) on a smooth geometrically connected curve. If the irreducible constituents \(\mathcal V_1,\dots,\mathcal V_n\) of \(\mathcal V\) are pure of weights \(w_1,\dots,w_n\), then \(\operatorname{dw}_\iota(\mathcal V)=\sum_a\operatorname{rank}(\mathcal V_a)w_a/r\). In particular a sheaf pure of weight \(w\) has determinant weight \(w\).
\end{bpstatement}



% object: sec-Determinant-weights (blueprint/chap-Purity-of-real-sheaves/sec-Determinant-weights/_section.md), content hidden
\subsection*{\texorpdfstring{14.2\quad Determinant weights}{14.2 Determinant weights}}
\phantomsection\label{bp31}%
\addcontentsline{toc}{subsection}{\texorpdfstring{14.2\quad Determinant weights}{14.2 Determinant weights}}
% text: prose-14-11
\phantomsection\label{bp32}%
Before purity is known, the only weight we can attach to an irreducible lisse Weil sheaf \(Q\) on a curve is its determinant weight \(\operatorname{dw}_\iota(Q)\), the weight of \(\det Q\) divided by the rank (Definition \hyperref[bp26]{14.7} and Corollary \hyperref[bp27]{14.8}). This section proves the two facts about determinant weights that Rankin's argument in Section \hyperref[bp64]{14.3} needs for the tensor and exterior powers of a real sheaf \(\mathcal G\). First, every simple factor of \(\mathcal G^{\otimes2m}\) has determinant weight at most \(2m\) times the largest determinant weight of the constituents of \(\mathcal G\), and this bounds the weights of \(H^2_c\). Second, the determinant weights of the simple factors of \(\bigwedge^r\mathcal G\) are known exactly, and this is what lets exterior powers separate the constituents.

The difficulty is that a determinant only sees an average. For irreducible \(A\) and \(B\) of ranks \(a\) and \(b\), we have \(\det(A\otimes B)=(\det A)^{\otimes b}\otimes(\det B)^{\otimes a}\). Hence the determinant weight of \(A\otimes B\) as a whole is \(\operatorname{dw}_\iota(A)+\operatorname{dw}_\iota(B)\). But \(A\otimes B\) is usually reducible, and which subquotients occur depends on the monodromy. The rank-weighted average of the determinant weights of its simple factors is the number above, but nothing so far prevents one factor from lying above it and another below. If the factors were pure, the determinant weight of each would be its weight, but purity is what we are trying to prove.

The model to keep in mind is a central character. If a group acts on \(V\) and some central element \(z\) acts on \(V\) by the scalar \(c\), then \(z\) acts on every subquotient of \(V^{\otimes a}\otimes(V^\vee)^{\otimes b}\) by \(c^{a-b}\), without any knowledge of how these representations decompose. We therefore look for a central element of positive degree in the monodromy group; its scalars will record weights.

Let \(H\subseteq G\) be the Zariski closures in \(\GL(V)\) of the images of \(\pi_1(U)\) and of \(W(U_0)\), the geometric and arithmetic monodromy groups. If \(H\) were trivial, the image \(\rho(s)\) of a degree-one lift would itself be central. In general we use the structure of the identity component,

\[
G^\circ=Z(G^\circ)^\circ H^\circ
\]

for \(V\) semisimple (Lemma \hyperref[bp49]{14.15}). The reason is that the Lie algebra of \(G^\circ\) is reductive, hence the sum of its center and its derived algebra, and since \(G/H\) is abelian, the derived algebra lies in the Lie algebra of \(H\). Accordingly, in Proposition \hyperref[bp54]{14.17} we write a power of \(\rho(s)\) as a central torus element times an element of \(H\), and then we take the product of its conjugates over \(G/G^\circ\). The result is an element \(z=h\rho(s)^N\), central in \(G\), with \(h\in H\) and \(N>0\).

On a simple factor \(Q\) of any tensor construction, \(z\) acts by a scalar \(c_Q\), and taking determinants gives \(c_Q^{\dim Q}=\det(h\mid Q)\det(s\mid Q)^N\). Since determinants have finite geometric monodromy, \(\det(h\mid Q)\) is a root of unity, and from this Lemma \hyperref[bp58]{14.18} concludes that \(\operatorname{dw}_\iota(Q)=\lambda_\iota(c_Q)/N\), where \(\lambda_\iota(a)=2\log_q|\iota(a)|\) is the additive function of Proposition \bpout{13.7}. Since scalars multiply under tensor products, we obtain Proposition \hyperref[bp61]{14.19}.

The facts on algebraic groups that we use, namely the relation between closed subgroups, their identity components and their Lie algebras, and the structure of a Lie algebra with a faithful semisimple representation, are proved in Appendix \bpout{F}.

% object: subsec-Representations-and-filtrations (blueprint/chap-Purity-of-real-sheaves/sec-Determinant-weights/subsec-Representations-and-filtrations/_section.md), content hidden
\subsubsection*{\texorpdfstring{14.2.1\quad Representations and filtrations}{14.2.1 Representations and filtrations}}
\phantomsection\label{bp33}%
\addcontentsline{toc}{subsubsection}{\texorpdfstring{14.2.1\quad Representations and filtrations}{14.2.1 Representations and filtrations}}
% text: prose-14-12
\phantomsection\label{bp34}%
We treat lisse sheaves as representations of the Weil group. From the geometry of the curve we use only one fact, proved in Section \hyperref[bp3]{14.1}: the geometric fundamental group acts through a finite group on every determinant. Everything else is linear algebra and the structure of algebraic groups.

Throughout, \(U_0\) is a smooth geometrically connected curve over \(k_0\) and \(s\in W(U_0)\) is a lift of degree one, as in Convention \bpslug{situation-{}\allowbreak{}weil}, so that \(W(U_0)=\pi_1(U)\rtimes\langle s\rangle\).

% object: definition-weil-representation (blueprint/chap-Purity-of-real-sheaves/sec-Determinant-weights/subsec-Representations-and-filtrations/definition-weil-representation/_section.md)
\phantomsection\label{bp35}%
\begin{bpstatement}{Definition~14.10}{ (Weil representation)}
A \emph{Weil representation} is a finite-dimensional \(\overline{\mathbf Q}_\ell\)-representation of \(W(U_0)\) whose restriction to \(\pi_1(U)\) is continuous.
\end{bpstatement}

% text: prose-14-13
\phantomsection\label{bp36}%
By Convention \bpslug{situation-{}\allowbreak{}weil} the Weil representations are the lisse Weil \(\overline{\mathbf Q}_\ell\)-sheaves on \(U_0\), and so, by Corollary \hyperref[bp27]{14.8}, \(\pi_1(U)\) acts through a finite group on the determinant of every Weil representation. A Weil representation \(V\) is semisimple if it is a direct sum of simple representations, and the semisimplification \(V^{\mathrm{ss}}\) of \(V\) is the direct sum of its composition factors.

Two elementary facts are needed. First, semisimplicity is preserved by restriction to a normal subgroup, which lets us apply Lie theory to the identity component of the monodromy group below. Second, the composition factors of tensor and exterior powers are computed from those of the given representations, so it suffices to prove the formulas for semisimple representations.

% object: lemma-restriction-normal-subgroup (blueprint/chap-Purity-of-real-sheaves/sec-Determinant-weights/subsec-Representations-and-filtrations/lemma-restriction-normal-subgroup/_section.md)
\phantomsection\label{bp37}%
\begin{bpstatementit}{Lemma~14.11}{ (Restriction to a normal subgroup)}
Let \(N\) be a normal subgroup of a group \(\Gamma\). The restriction to \(N\) of a finite-dimensional semisimple \(\overline{\mathbf Q}_\ell\)-representation of \(\Gamma\) is semisimple.
\end{bpstatementit}

% proof: proof/lemma-restriction-normal-subgroup
\begin{proof}[{Proof}]\phantomsection\label{bp38}%
It suffices to consider a simple \(\Gamma\)-module \(V\). A nonzero \(N\)-submodule \(A\) of least dimension is simple, and since \(N\) is normal, each translate \(gA\) is again a simple \(N\)-submodule. The sum of all these translates is a nonzero \(\Gamma\)-submodule, hence equals \(V\). It remains to see that a finite sum of simple submodules is semisimple. Indeed, the intersection of a simple submodule with the sum of those already chosen is either zero or the whole submodule, so we add it in the first case and omit it in the second, and after finitely many steps we obtain a direct-sum decomposition.
\end{proof}

% text: prose-14-14
\phantomsection\label{bp39}%
The second fact concerns representations that are not semisimple. A nonsplit extension of \(\overline{\mathbf Q}_\ell(-1)\) by \(\overline{\mathbf Q}_\ell\), such as the twist \(\mathcal E(-1)\) of the Kummer extension of Example \bpout{13.4}, does not decompose. Yet the composition factors of its tensor and exterior powers are determined by those of the sheaf, as the next lemma shows by filtering. This is what lets us pass to semisimplifications in Proposition \hyperref[bp61]{14.19}.

% object: lemma-tensor-exterior-powers-filtered-representations (blueprint/chap-Purity-of-real-sheaves/sec-Determinant-weights/subsec-Representations-and-filtrations/lemma-tensor-exterior-powers-filtered-representations/_section.md)
\phantomsection\label{bp40}%
\begin{bpstatementit}{Lemma~14.12}{ (Tensor and exterior powers of filtered representations)}
Suppose \(V\) and \(V'\) have composition factors \(A_1,\ldots,A_t\) and \(B_1,\ldots,B_u\). There is a \(W(U_0)\)-stable filtration of \(V\otimes V'\) with graded pieces \(A_a\otimes B_b\), one for each pair \((a,b)\). There is a \(W(U_0)\)-stable filtration of \(\bigwedge^rV\) with graded pieces

\[
\bigotimes_{a=1}^t\bigwedge^{n_a}A_a,
\qquad 0\le n_a\le\dim A_a,\qquad\sum_an_a=r.
\]

The simple factors of \(V\otimes V'\) and \(\bigwedge^rV\), counted with multiplicity, are those of the respective graded pieces. Every graded piece of \(\bigwedge^rV\) displayed above is nonzero.
\end{bpstatementit}

% proof: proof/lemma-tensor-exterior-powers-filtered-representations
\begin{proof}[{Proof}]\phantomsection\label{bp41}%
We construct the two filtrations and then compare simple factors by the Jordan–Hölder argument. The tensor product is filtered by the filtration on \(V\), and each quotient \(A_a\otimes V'\) by the filtration on \(V'\); since tensoring over a field is exact, the quotients are as asserted.

For the exterior power we cannot simply filter factor by factor, because a wedge product may use several vectors from the same composition factor. Instead we sort wedges of basis vectors by how many vectors they take from each factor, after lifting bases of the \(A_a\) to a basis of \(V\) adapted to its filtration. Thus a wedge of basis vectors determines a tuple \((n_a)\), by counting how many vectors it uses from each lifted basis. Since the basis is adapted, the group \(W(U_0)\) carries a lifted basis vector of \(A_a\) into the span of the lifted bases of \(A_1,\ldots,A_a\). Hence the action is block upper triangular: each term of the image of a basis wedge either has the same tuple or has a tuple obtained from it by moving some of its vectors to smaller indices \(a\). Choose a total order on the tuples compatible with these moves. Then the spans of the initial segments of basis wedges form a stable filtration, and on the quotient indexed by \((n_a)\) the diagonal blocks act as on \(\bigotimes_a\bigwedge^{n_a}A_a\), a space of dimension \(\prod_a\binom{\dim A_a}{n_a}>0\). For example, for \(V\) an extension of \(A_2\) by \(A_1\), both of dimension one, and \(r=1\), the tuples are \((1,0)\) and \((0,1)\), and the filtration is the given one. For \(r=2\) there is the single tuple \((1,1)\), and \(\bigwedge^2V=A_1\otimes A_2\).

The graded pieces of the filtrations just built need not be simple, so to compare simple factors we refine. By the Schreier refinement theorem, any two finite filtrations have refinements with isomorphic graded pieces, matched by the Zassenhaus isomorphism \((V_{i-1}+V_i\cap W_j)/(V_{i-1}+V_i\cap W_{j-1})\simeq
(W_{j-1}+W_j\cap V_i)/(W_{j-1}+W_j\cap V_{i-1})\); for two composition series this is the Jordan–Hölder theorem. Applying the same refinement to a composition series and to the tensor or exterior filtration, we obtain the statement about simple factors. In particular, these multisets do not change when \(V\) and \(V'\) are replaced by their semisimplifications.
\end{proof}


% object: subsec-Central-elements (blueprint/chap-Purity-of-real-sheaves/sec-Determinant-weights/subsec-Central-elements/_section.md), content hidden
\subsubsection*{\texorpdfstring{14.2.2\quad Central elements}{14.2.2 Central elements}}
\phantomsection\label{bp42}%
\addcontentsline{toc}{subsubsection}{\texorpdfstring{14.2.2\quad Central elements}{14.2.2 Central elements}}
% text: prose-14-15
\phantomsection\label{bp43}%
In the arithmetic monodromy group we construct a central element of positive degree (Proposition \hyperref[bp54]{14.17}). It acts by a scalar on each simple factor, and this scalar determines the determinant weight of the factor. As explained above, we construct the element for semisimple representations, and then use stable filtrations to obtain the formulas for arbitrary representations.

The degree must be positive. Indeed, an element of \(H\) that is central in \(G\) also acts by a scalar on each simple factor \(Q\). But, as for every element of \(H\), its determinant on \(Q\) is a root of unity, so under \(\iota\) the scalar has absolute value one, and it records no weight. By contrast, an element \(z=h\rho(s)^N\) with \(N>0\) contains \(N\) factors of Frobenius, and its scalar \(c_Q\) on \(Q\) satisfies \(\lambda_\iota(c_Q)=N\operatorname{dw}_\iota(Q)\) (Lemma \hyperref[bp58]{14.18}).

We produce the central element by Lie theory, which applies to algebraic groups and not to the abstract images of \(\pi_1(U)\) and \(W(U_0)\). For \(\mathcal B_b\) with \(b\) not a root of unity, the image of \(W(U_0)\) in \(\GL_1\) is the infinite cyclic group generated by \(b\), which is not an algebraic subgroup; its Zariski closure is all of \(\GL_1\), since the proper closed subgroups of \(\mathbf G_m\) are finite. Passing to the closure loses nothing that we need. Indeed, stabilizers are closed, so a subspace of \(V\) is stable under a group of matrices if and only if it is stable under its Zariski closure. For this reason Lemma \hyperref[bp49]{14.15} and Proposition \hyperref[bp54]{14.17} are statements about these closures.

% object: definition-monodromy-groups (blueprint/chap-Purity-of-real-sheaves/sec-Determinant-weights/subsec-Central-elements/definition-monodromy-groups/_section.md)
\phantomsection\label{bp44}%
\begin{bpstatement}{Definition~14.13}{ (Monodromy groups)}
Let \(\rho:W(U_0)\to\GL(V)\) be a Weil representation. Its \emph{geometric} and \emph{arithmetic monodromy groups} are the Zariski closures \(H\) of \(\rho(\pi_1(U))\) and \(G\) of \(\rho(W(U_0))\) in \(\GL(V)\).
\end{bpstatement}

% text: prose-14-16
\phantomsection\label{bp45}%
We keep the letters \(H\) and \(G\) for these two groups in the rest of this section.

% object: lemma-normality-geometric-monodromy (blueprint/chap-Purity-of-real-sheaves/sec-Determinant-weights/subsec-Central-elements/lemma-normality-geometric-monodromy/_section.md)
\phantomsection\label{bp46}%
\begin{bpstatementit}{Lemma~14.14}{ (Normality of geometric monodromy)}
The groups \(H\) and \(G\) are reduced closed algebraic subgroups of \(\GL(V)\), and \(H\triangleleft G\). Moreover every commutator of two elements of \(G\) lies in \(H\).
\end{bpstatementit}

% proof: proof/lemma-normality-geometric-monodromy
\begin{proof}[{Proof}]\phantomsection\label{bp47}%
Each assertion holds for the images of \(\pi_1(U)\) and \(W(U_0)\), and we pass to Zariski closures by continuity. First, multiplication and inversion are continuous for the Zariski topology, so the closure of a subgroup is a subgroup. Next, conjugation by \(\rho(W(U_0))\) preserves \(H\) because \(\pi_1(U)\) is normal in \(W(U_0)\), and since the normalizer of a closed subgroup is closed, conjugation by \(G\) preserves \(H\) as well. Finally, since \(W(U_0)/\pi_1(U)\simeq\mathbf Z\) is abelian, commutators of elements of \(\rho(W(U_0))\) lie in \(\rho(\pi_1(U))\). As the commutator map \(G\times G\to\GL(V)\) is continuous, the same holds on the closures.
\end{proof}

% text: prose-14-17
\phantomsection\label{bp48}%
For example, for \(V=\mathcal B_b\) of Proposition \bpout{13.7} the group \(H\) is trivial. Hence \(\rho(s)=b\) is central in \(G\), and the determinant weight of \(\mathcal B_b\) is \(\lambda_\iota(b)\).

% object: lemma-identity-component (blueprint/chap-Purity-of-real-sheaves/sec-Determinant-weights/subsec-Central-elements/lemma-identity-component/_section.md)
\phantomsection\label{bp49}%
\begin{bpstatementit}{Lemma~14.15}{ (The identity component)}
If \(V\) is a semisimple \(W(U_0)\)-representation, then

\[
G^\circ=Z(G^\circ)^\circ H^\circ.
\]
\end{bpstatementit}

% proof: proof/lemma-identity-component
\begin{proof}[{Proof}]\phantomsection\label{bp50}%
A connected algebraic group is determined by its Lie algebra among its closed connected subgroups, so it suffices to show that the Lie algebras of \(Z(G^\circ)^\circ\) and \(H^\circ\) together span that of \(G^\circ\). We show that this Lie algebra is reductive, hence the sum of its center and its derived algebra, and that the derived algebra lies in the Lie algebra of \(H\) because commutators lie in \(H\).

\emph{Reductivity.} Write \(\mathfrak g=\operatorname{Lie}(G^\circ)\). Then

\[
\mathfrak g=Z(\mathfrak g)+[\mathfrak g,\mathfrak g],
\]

by Proposition \bpout{F.8}, once we know that \(V\) is a faithful semisimple \(\mathfrak g\)-module. The \(\rho(W(U_0))\)-stable subspaces are precisely the \(G\)-stable subspaces, because the stabilizer of a subspace is closed. Hence \(V\) is semisimple as a \(G(\overline{\mathbf Q}_\ell)\)-module, and by Lemma \hyperref[bp37]{14.11} it is also semisimple as a \(G^\circ(\overline{\mathbf Q}_\ell)\)-module. Then Lemma \bpout{F.3} makes it a faithful semisimple \(\mathfrak g\)-module.

\emph{The center and the derived algebra.} Put \(T=Z(G^\circ)^\circ\) and \(\mathfrak h=\operatorname{Lie}(H^\circ)\). The comparison between groups and Lie algebras in Lemma \bpout{F.3} gives \(\operatorname{Lie}(T)=Z(\mathfrak g)\). Moreover \([\mathfrak g,\mathfrak g]\subseteq\mathfrak h\), because the commutator morphism \(G\times G\to G\) factors through the closed subgroup \(H\): it does so on \(\overline{\mathbf Q}_\ell\)-points by Lemma \hyperref[bp46]{14.14}, and \(G\times G\) is reduced. To deduce the inclusion, let \(X,Y\in\mathfrak g\) and put \(R=\overline{\mathbf Q}_\ell[\varepsilon,\eta]/(\varepsilon^2,\eta^2)\). The commutator of \(1+\varepsilon X\) and \(1+\eta Y\) in \(G(R)\) is \(1+\varepsilon\eta[X,Y]\), and it lies in \(H(R)\). Here the subring \(\overline{\mathbf Q}_\ell[\varepsilon\eta]\subset R\) is a ring of dual numbers, and the equations of \(H\) vanish on \(1+\varepsilon\eta[X,Y]\) already over this subring. Hence \([X,Y]\in\operatorname{Lie}(H)=\mathfrak h\), and

\[
\mathfrak g=\operatorname{Lie}(T)+\mathfrak h.
\]

\emph{Back to groups.} The multiplication \(T\times H^\circ\to G^\circ\) is a homomorphism, because \(T\) is central. By Lemma \bpout{F.2} its image is a closed connected subgroup, whose Lie algebra contains both summands on the right, hence equals that of \(G^\circ\). By the same lemma the image is therefore \(G^\circ\).
\end{proof}

% text: prose-14-18
\phantomsection\label{bp51}%
Rankin's argument needs the determinant weights of the simple factors of \(\mathcal G^{\otimes2m}\), of the exterior powers of \(\mathcal G\) and of the dual, so we give a name to the representations built in this way; the monodromy group acts on all of them.

% object: definition-tensor-construction (blueprint/chap-Purity-of-real-sheaves/sec-Determinant-weights/subsec-Central-elements/definition-tensor-construction/_section.md)
\phantomsection\label{bp52}%
\begin{bpstatement}{Definition~14.16}{ (Tensor construction)}
A \emph{tensor construction} on a Weil representation \(V\) is a representation obtained from \(V\) and \(V^\vee\) by finitely many tensor products, direct sums, exterior powers and subquotients by \(W(U_0)\)-stable subspaces.
\end{bpstatement}

% text: prose-14-19
\phantomsection\label{bp53}%
Since the stabilizer of a subspace is closed, a \(W(U_0)\)-stable subspace is \(G\)-stable, and so \(G\) acts on every tensor construction through its action on \(V\).

% object: proposition-central-element-arithmetic-monodromy-group (blueprint/chap-Purity-of-real-sheaves/sec-Determinant-weights/subsec-Central-elements/proposition-central-element-arithmetic-monodromy-group/_section.md)
\phantomsection\label{bp54}%
\begin{bpstatementit}{Proposition~14.17}{ (A central element of the arithmetic monodromy group)}
Let \(V_1,\ldots,V_n\) be semisimple Weil representations and form the monodromy groups \(H\subseteq G\) of \(V=\bigoplus_iV_i\), and let \(s\in W(U_0)\) have degree one. There are an integer \(N>0\) and \(h\in H(\overline{\mathbf Q}_\ell)\) such that

\[
z=h\,\rho(s)^N\in Z(G)(\overline{\mathbf Q}_\ell).
\]

The action of \(z\) on each tensor construction on \(V\) preserves every \(W(U_0)\)-stable subquotient and is scalar on each simple subquotient.
\end{bpstatementit}

% proof: proof/proposition-central-element-arithmetic-monodromy-group
\begin{proof}[{Proof}]\phantomsection\label{bp55}%
We write a power of \(\rho(s)\) as a central torus element times an element of \(H\), by Lemma \hyperref[bp49]{14.15}, and average the torus part over the components of \(G\) so that it becomes central in all of \(G\). The factors from \(H\) that accumulate in this process are collected in \(h\); their determinants are roots of unity, so by Lemma \hyperref[bp58]{14.18} they do no harm. Note that the direct sum \(V\) of the \(V_i\) is again semisimple, so Lemma \hyperref[bp49]{14.15} applies to it.

Since the component group of \(G\) is finite (Lemma \bpout{F.2}), we may choose \(m>0\) with \(\rho(s)^m\in G^\circ\). By Lemma \hyperref[bp49]{14.15}, we can write \(\rho(s)^m=ct\) with \(c\in Z(G^\circ)^\circ\) and \(t\in H^\circ\). Choose representatives \(g_1,\ldots,g_b\) of the components of \(G\) and set

\[
z=\prod_{j=1}^bg_jcg_j^{-1},\qquad N=mb.
\]

The element \(c\) is central in \(G^\circ\) but perhaps not in \(G\), since the other components may conjugate it nontrivially. The product \(z\) of its conjugates, one for each component, is however central in \(G\). Indeed, the factors lie in the abelian group \(Z(G^\circ)^\circ\), and since replacing \(g_j\) by \(g_ju\) with \(u\in G^\circ\) does not change the corresponding factor, conjugation by an element of \(G\) permutes the factors. Next, every commutator lies in \(H\), so each factor has the same image as \(c\) in \(G/H\). Since \(c\in H\rho(s)^m\), we get \(z\in H\rho(s)^{mb}\), which gives \(h\) and \(N=mb>0\).

Being central in \(G\), the element \(z\) commutes with the action of \(W(U_0)\) on tensor constructions. Moreover, as noted after the definition of tensor constructions, every \(W(U_0)\)-stable subspace of a tensor construction is \(G\)-stable. Hence \(z\) acts on the \(W(U_0)\)-stable subquotients and commutes with the action of \(W(U_0)\) on them. On a simple subquotient, therefore, an eigenspace of \(z\) is nonzero and \(W(U_0)\)-stable, hence is the whole space; in other words, \(z\) acts by a scalar.
\end{proof}


% object: subsec-Tensor-operations (blueprint/chap-Purity-of-real-sheaves/sec-Determinant-weights/subsec-Tensor-operations/_section.md), content hidden
\subsubsection*{\texorpdfstring{14.2.3\quad Tensor operations}{14.2.3 Tensor operations}}
\phantomsection\label{bp56}%
\addcontentsline{toc}{subsubsection}{\texorpdfstring{14.2.3\quad Tensor operations}{14.2.3 Tensor operations}}
% text: prose-14-20
\phantomsection\label{bp57}%
The central element acts on each simple factor of a tensor construction by a scalar, and from these scalars we read off determinant weights. Since Weil representations are the lisse Weil \(\overline{\mathbf Q}_\ell\)-sheaves, Definition \hyperref[bp26]{14.7} applies to them.

% object: lemma-determinant-weight-from-central-scalar (blueprint/chap-Purity-of-real-sheaves/sec-Determinant-weights/subsec-Tensor-operations/lemma-determinant-weight-from-central-scalar/_section.md)
\phantomsection\label{bp58}%
\begin{bpstatementit}{Lemma~14.18}{ (Determinant weight from a central scalar)}
In the setting of Proposition \hyperref[bp54]{14.17}, let \(Q\) be a nonzero subquotient of a tensor construction on \(V\) on which the element \(z\) acts by a scalar \(c_Q\). Then

\[
\operatorname{dw}_\iota(Q)=\frac{\lambda_\iota(c_Q)}{N}.
\]
\end{bpstatementit}

% proof: proof/lemma-determinant-weight-from-central-scalar
\begin{proof}[{Proof}]\phantomsection\label{bp59}%
The determinant of \(Q\) is an algebraic character of \(H\), and its image on the dense subgroup \(\rho(\pi_1(U))\) is finite, and hence its image on all of \(H\) lies in the same finite group of roots of unity. Taking determinants in \(z=h\rho(s)^N\), we find

\[
c_Q^{\dim Q}=\det(h\mid Q)\det(s\mid Q)^N.
\]

Now apply \(\lambda_\iota\), which vanishes on the first factor, and divide by \(N\dim Q\).
\end{proof}

% text: prose-14-21
\phantomsection\label{bp60}%
The lemma reduces determinant weights of tensor constructions to a computation with the scalars on the composition factors of \(V\). By Proposition \hyperref[bp54]{14.17}, the element \(z\) acts by a scalar on every simple subquotient of a tensor construction on \(V\), so the formula of the lemma applies to each of them. Since the scalars by which a single element acts multiply under tensor products and exterior powers, this gives the following formulas.

% object: proposition-determinant-weights-tensor-operations (blueprint/chap-Purity-of-real-sheaves/sec-Determinant-weights/subsec-Tensor-operations/proposition-determinant-weights-tensor-operations/_section.md)
\phantomsection\label{bp61}%
\begin{bpstatementit}{Proposition~14.19}{ (Determinant weights of tensor operations)}
Let \(V\) and \(V'\) be nonzero Weil representations, let \(A_1,\ldots,A_t\) be the composition factors of \(V\), and put \(\gamma_a=\operatorname{dw}_\iota(A_a)\).

\begin{enumerate}
\item \(\operatorname{DW}(V\otimes V')=\operatorname{DW}(V)+\operatorname{DW}(V')\).
\item \(\operatorname{DW}(\bigwedge^rV)=\bigl\{\sum_an_a\gamma_a :
      0\le n_a\le\dim A_a,\ \sum_an_a=r\bigr\}\) for \(0\le r\le\dim V\).
\item \(\operatorname{DW}(V^\vee)=-\operatorname{DW}(V)\).
\end{enumerate}
\end{bpstatementit}

% proof: proof/proposition-determinant-weights-tensor-operations
\begin{proof}[{Proof}]\phantomsection\label{bp62}%
We use one central element for \(V\) and \(V'\) simultaneously, so that its scalars on the factors of \(V\otimes V'\) are products of its scalars on the factors of \(V\) and of \(V'\). Since Proposition \hyperref[bp54]{14.17} requires semisimple representations, we apply it to \(V^{\mathrm{ss}}\oplus V'^{\mathrm{ss}}\). This changes nothing, because by Lemma \hyperref[bp40]{14.12} the simple factors of \(V\otimes V'\) and \(\bigwedge^rV\) depend only on the composition factors of \(V\) and \(V'\). Now let \(c_a\) and \(c'_b\) be the scalars by which its central element of degree \(N\) acts on the composition factors \(A_a\) of \(V\) and \(A'_b\) of \(V'\). By Lemma \hyperref[bp58]{14.18}, \(\gamma_a=\lambda_\iota(c_a)/N\) and \(\operatorname{dw}_\iota(A'_b)=\lambda_\iota(c'_b)/N\).

(1) The central element acts by \(c_ac'_b\) on \(A_a\otimes A'_b\), and hence on every simple factor of it. By the same lemma, each such factor has determinant weight \(\gamma_a+\operatorname{dw}_\iota(A'_b)\). Moreover, each \(A_a\otimes A'_b\) is nonzero, and Lemma \hyperref[bp40]{14.12} identifies the simple factors of these tensor products with those of \(V\otimes V'\).

(2) The central element acts on \(\bigotimes_a\bigwedge^{n_a}A_a\) by \(\prod_ac_a^{n_a}\). So each simple factor of this nonzero representation has determinant weight \(\sum_an_a\gamma_a\), and Lemma \hyperref[bp40]{14.12} identifies these factors with those of \(\bigwedge^rV\).

(3) We have \(\det(s\mid A^\vee)=\det(s\mid A)^{-1}\), and duality carries a composition series of \(V\) to one of \(V^\vee\) whose factors are the duals \(A_a^\vee\).
\end{proof}

% text: prose-14-22
\phantomsection\label{bp63}%
In particular every simple factor of \(V^{\otimes m}\) has determinant weight at most \(m\max\operatorname{DW}(V)\). In Section \hyperref[bp80]{14.4} we use this bound, and part (2), for the constituents of a real sheaf before we know that they are pure.



% object: sec-Rankins-bound (blueprint/chap-Purity-of-real-sheaves/sec-Rankins-bound/_section.md), content hidden
\subsection*{\texorpdfstring{14.3\quad Rankin's bound}{14.3 Rankin's bound}}
\phantomsection\label{bp64}%
\addcontentsline{toc}{subsection}{\texorpdfstring{14.3\quad Rankin's bound}{14.3 Rankin's bound}}
% text: prose-14-23
\phantomsection\label{bp65}%
We bound the local eigenvalues of a real sheaf by the eigenvalues of Frobenius on \(H^2_c\) of its even tensor powers (Proposition \hyperref[bp78]{14.25}). Only \(H^2_c\) is needed, and Lemma \hyperref[bp66]{14.20} computes it by Poincaré duality as the coinvariants of the geometric monodromy.

Rankin's argument uses one fact from each side of the trace formula. On the cohomological side, an affine curve has \(H^0_c=0\), so only \(H^2_c\) can produce poles of the \(L\)-function, and we can compute this group without knowing anything about \(H^1_c\): its eigenvalues are \(q\) times eigenvalues of Frobenius on a quotient of the generic fibre. On the side of the Euler product, if the local factors have nonnegative coefficients, then each of them is holomorphic wherever the product is.

For example, for the constant sheaf on \(\mathbf G_m\) we have

\[
L(\mathbf G_m,\overline{\mathbf Q}_\ell,t)=\prod_{x\in|\mathbf G_m|}(1-t^{\deg(x)})^{-1}
=\frac{1-t}{1-qt},
\]

where \(1-t\) comes from \(H^1_c\) and the pole at \(t=q^{-1}\) from \(H^2_c=\overline{\mathbf Q}_\ell(-1)\). Each local factor has its poles on the circle \(|t|=1\), outside the disc \(|t|<q^{-1}\), as Lemma \hyperref[bp75]{14.24} predicts. In general we apply the two facts to the even tensor powers of a real sheaf, whose \(H^2_c\) is bounded by the determinant weights of Section \hyperref[bp31]{14.2}, and this gives Proposition \hyperref[bp78]{14.25}.

% object: lemma-degree-zero-top-cohomology (blueprint/chap-Purity-of-real-sheaves/sec-Rankins-bound/lemma-degree-zero-top-cohomology/_section.md)
\phantomsection\label{bp66}%
\begin{bpstatementit}{Lemma~14.20}{ (Degree zero and top cohomology)}
Let \(U_0/\mathbf F_q\) be a smooth geometrically connected curve, let \(\mathcal F\) be a lisse Weil sheaf on \(U_0\), and let \(V\) be its geometric generic fibre, a representation of the geometric fundamental group \(\pi_1(U)\).  There are canonical Frobenius-equivariant identifications

\[
H^0(U,\mathcal F)=V^{\pi_1(U)},\qquad
H_c^2(U,\mathcal F)=V_{\pi_1(U)}(-1).
\]

There is no compactly supported cohomology outside degrees \(0,1,2\). If \(U_0\) is affine, \(H_c^0(U,\mathcal F)=0\). If \(\mathcal F\) is punctually of weights \(\le a\), the weights of \(H^0\) are at most \(a\) and those of \(H_c^2\) at most \(a+2\).  If \(\mathcal F\) is pure of weight \(a\), the two groups, when nonzero, are pure of weights \(a\) and \(a+2\), respectively.
\end{bpstatementit}

% proof: proof/lemma-degree-zero-top-cohomology
\begin{proof}[{Proof}]\phantomsection\label{bp67}%
We deduce both formulas from the description of a lisse sheaf by its representation, using Poincaré duality to reduce \(H^2_c\) to \(H^0\), and then we read off the weights from the stalks.

A global section of a local system is determined by its value at \(\bar\eta\), and a vector of the fibre extends to a global section exactly when it is fixed by the monodromy. This is the first formula.  For the second, by degreewise Poincaré duality (Corollary \bpout{10.13}) we have

\[
H_c^2(U,\mathcal F)^\vee
=H^0(U,\mathcal F^\vee(1))=(V^\vee)^{\pi_1(U)}(1).
\]

Since the dual of \((V^\vee)^{\pi_1(U)}\) is \(V_{\pi_1(U)}\), this gives the second formula, with its twist. Moreover, the identification is induced by the canonical duality pairing and is therefore Frobenius equivariant.

Duality also gives the vanishing outside degrees \(0,1,2\). Indeed, for \(i<0\) there is nothing to prove, and for \(i>2\) the space \(H_c^i(U,\mathcal F)\) is dual to \(H^{2-i}(U,\mathcal F^\vee(1))\), which is zero because \(2-i<0\). For torsion coefficients the vanishing comes instead from the bound on cohomological dimension of Appendix \bpout{D}.

Suppose that \(U_0\) is affine, and choose a smooth proper completion of \(U\). Its boundary is nonempty, because a connected proper curve has only constant global functions and hence is not affine. A compactly supported section extends by zero to the completion, so it vanishes on a neighbourhood of the boundary. Since a section of a lisse sheaf on a connected curve is determined by its generic value, the section is zero, and \(H_c^0(U,\mathcal F)=0\).

For the weights we compare the global Frobenius on these two spaces with the local Frobenius at a closed point. Since \(\pi_1(U)\) acts trivially on \(V^{\pi_1(U)}\) and on \(V_{\pi_1(U)}\), the action of \(W(U_0)\) on them factors through the degree quotient.  If \(\beta\) is an eigenvalue of the operator of degree one, then \(\beta^{\deg(x)}\) is an eigenvalue of \(\operatorname{Frob}_x\) on \(\mathcal F_{\bar x}\) for every closed point \(x\). Indeed, \(V^{\pi_1(U)}\) is a subspace and \(V_{\pi_1(U)}\) a quotient of \(V\simeq\mathcal F_{\bar x}\), both stable under \(\operatorname{Frob}_x\in\pi_1(U)s^{\deg(x)}\), which acts on them as \(s^{\deg(x)}\).  If \(\mathcal F\) is punctually of weights \(\le a\), it follows that \(|\iota(\beta)|\le q^{a/2}\), and if \(\mathcal F\) is pure of weight \(a\), equality holds. Finally, the twist \((-1)\) multiplies each global eigenvalue by \(q\), which gives the stated weights for \(H_c^2\).
\end{proof}

% text: prose-14-24
\phantomsection\label{bp68}%
For the constant sheaf on \(\mathbf G_m\) the lemma gives \(H^2_c=\overline{\mathbf Q}_\ell(-1)\), of weight \(2\), as in the computation at the start of the section. Now take a Kummer sheaf of a nontrivial character on \(\mathbf G_m\). Then \(H^0_c=0\) because \(\mathbf G_m\) is affine, and \(H^2_c=0\) because the coinvariants vanish.

The positivity we want is \(a^{2m}\ge0\) for the traces \(a\) of all powers \(\operatorname{Frob}_x^n\) on \(\mathcal G\), and for this all these traces must be real. It is not enough to ask that the trace of \(\operatorname{Frob}_x\) itself be real. For example, if the eigenvalues at \(x\) are \(i\) and \(1-i\), the trace is \(1\), but the trace of the square is \(-1+(1-i)^2=-1-2i\). By the Newton identities, on the other hand, all the power sums are real exactly when the characteristic polynomial is real, and this is the condition below. Its basic examples are the constant sheaf, its Tate twists and Kummer sheaves of quadratic characters. More generally, in Chapter \bpout{15}, Proposition \bpout{15.9} makes every pure sheaf a summand of a real one after a twist. The condition enters the proof of Theorem \hyperref[bp82]{14.26} through Lemma \hyperref[bp71]{14.22}.

% object: definition-real-sheaves (blueprint/chap-Purity-of-real-sheaves/sec-Rankins-bound/definition-real-sheaves/_section.md)
\phantomsection\label{bp69}%
\begin{bpstatement}{Definition~14.21}{ (Real sheaves)}
A Weil sheaf \(\mathcal G\) is \(\iota\)-real if \(\iota\det(1-T\operatorname{Frob}_x\mid\mathcal G_{\bar x})\) has real coefficients for every closed point; equivalently, the multiset of its complex eigenvalues at each point is stable under complex conjugation.  In particular, the traces of all positive powers of local Frobenius are real.
\end{bpstatement}

% text: prose-14-25
\phantomsection\label{bp70}%
The next lemma is the \(\ell\)-adic form of Rankin's trick. Since the even powers of a real trace are nonnegative, the Euler factors of \(\mathcal G^{\otimes2m}\) have nonnegative coefficients.

% object: lemma-positivity-euler-factors (blueprint/chap-Purity-of-real-sheaves/sec-Rankins-bound/lemma-positivity-euler-factors/_section.md)
\phantomsection\label{bp71}%
\begin{bpstatementit}{Lemma~14.22}{ (Positivity of Euler factors)}
Let \(\mathcal G\) be an \(\iota\)-real Weil sheaf, let \(m\ge1\), and let \(x\) be a closed point. Then the power series \(\iota\det(1-\operatorname{Frob}_xt^{\deg(x)}\mid\mathcal G^{\otimes2m}_{\bar x})^{-1}\) has nonnegative real coefficients and constant coefficient one.
\end{bpstatementit}

% proof: proof/lemma-positivity-euler-factors
\begin{proof}[{Proof}]\phantomsection\label{bp72}%
We write an inverse characteristic polynomial as the exponential of a series in the traces of powers, whose coefficients are then visibly nonnegative. For a square matrix \(A\), triangularizing \(A\) and expanding the logarithm of each factor we get the formal identity

\[
\det(1-tA)^{-1}
=\exp\left(\sum_{n\ge1}\operatorname{Tr}(A^n)\frac{t^n}{n}\right).
\]

On an even tensor power, the trace of the \(n\)th power of Frobenius is \(\operatorname{Tr}(\operatorname{Frob}_x^n\mid\mathcal G)^{2m}\). Since \(\mathcal G\) is real, this is nonnegative after applying \(\iota\).  Every coefficient of the exponential is then a finite sum of products of nonnegative numbers. This remains true after replacing \(t\) by \(t^{\deg(x)}\).
\end{proof}

% object: example-even-tensor-powers (blueprint/chap-Purity-of-real-sheaves/sec-Rankins-bound/example-even-tensor-powers/_section.md)
\phantomsection\label{bp73}%
\begin{bpstatement}{Example~14.23}{ (Why even tensor powers)}
The sheaf \(\mathcal B_{-1}\) is real, since \(\operatorname{Frob}_x\) acts by \((-1)^{\deg(x)}\), but its Euler factor at a point of degree one is \((1+t)^{-1}=1-t+t^2-\cdots\), with negative coefficients. Its square \(\mathcal B_{-1}\otimes\mathcal B_{-1}\) is the constant sheaf, whose Euler factors \((1-t^{\deg(x)})^{-1}\) have nonnegative coefficients, as the lemma predicts for \(m=1\).
\end{bpstatement}

% text: prose-14-26
\phantomsection\label{bp74}%
Without positivity a factor can converge on a smaller disc than the product, as in the cancellation at the start of this chapter. With positivity, however, the coefficients of each factor are bounded by those of the product.

% object: lemma-positivity-convergence (blueprint/chap-Purity-of-real-sheaves/sec-Rankins-bound/lemma-positivity-convergence/_section.md)
\phantomsection\label{bp75}%
\begin{bpstatementit}{Lemma~14.24}{ (Positivity and convergence)}
Let

\[
L(t)=\prod_j f_j(t),\qquad f_j\in1+t\mathbf R_{\ge0}[[t]],
\]

be a coefficientwise locally finite product. If \(L\) is the Taylor series of a function holomorphic on \(|t|<R\), then every \(f_j\) has radius of convergence at least \(R\), and if \(f_j=P_j^{-1}\) with \(P_j(0)=1\), then \(P_j\) has no zero in that disc.
\end{bpstatementit}

% proof: proof/lemma-positivity-convergence
\begin{proof}[{Proof}]\phantomsection\label{bp76}%
We compare coefficients and then apply Cauchy's estimate. To compare, write \(f_j=\sum_na_{j,n}t^n\) and \(L=\sum_nA_nt^n\). Since the product is coefficientwise locally finite, \(A_n\) is already the coefficient of \(t^n\) in a finite subproduct containing \(f_j\), which is \(f_j\) times a series \(g\) with constant coefficient one and nonnegative coefficients. The coefficient of \(t^n\) in \(f_jg\) is \(a_{j,n}\) plus a sum of products of nonnegative numbers, so that \(0\le a_{j,n}\le A_n\); positivity enters only here.  Next, for every \(0<r<R\), Cauchy's estimate gives \(|A_n|\le C_rr^{-n}\). Hence the radius of convergence of \(f_j\) is at least \(r\) for every such \(r\), and so at least \(R\).  A zero of \(P_j\) in the disc would be a pole of \(f_j=1/P_j\), which is impossible.
\end{proof}

% text: prose-14-27
\phantomsection\label{bp77}%
In the proposition below, the twist \((-1)\) on \(H^2_c\) costs a factor \(q\), and this loss disappears when we take \(2m\)th roots and let \(m\) grow.

% object: proposition-bound-largest-determinant-weight (blueprint/chap-Purity-of-real-sheaves/sec-Rankins-bound/proposition-bound-largest-determinant-weight/_section.md)
\phantomsection\label{bp78}%
\begin{bpstatementit}{Proposition~14.25}{ (A bound by the largest determinant weight)}
Let \(U_0/\mathbf F_q\) be a smooth geometrically connected affine curve and let \(\mathcal G\ne0\) be a lisse \(\iota\)-real Weil sheaf.  If \(\beta\) is the largest determinant weight of its irreducible Weil constituents, then

\[
|\iota(\alpha)|^2\le q_x^\beta
\quad(x\in|U_0|,\ \alpha\text{ an eigenvalue on }\mathcal G_{\bar x}).
\]
\end{bpstatementit}

% proof: proof/proposition-bound-largest-determinant-weight
\begin{proof}[{Proof}]\phantomsection\label{bp79}%
We bound the poles of \(L(U_0,\mathcal G^{\otimes2m},t)\) by determinant weights through \(H^2_c\), and positivity then passes the bound to each local factor.

We may assume that \(\mathcal G\) is semisimple. Indeed, by Lemma \bpout{13.5}, replacing \(\mathcal G\) by its semisimplification changes neither its local characteristic polynomials nor the determinant weights of its constituents.

\emph{The top cohomology.} Fix \(m\ge1\) and put \(\mathcal T=\mathcal G^{\otimes2m}\). We show that every eigenvalue \(\mu\) of \(F\) on \(H_c^2(U,\mathcal T)\) satisfies \(|\iota(\mu)|\le q^{m\beta+1}\); these eigenvalues determine the possible poles of the \(L\)-function. First, every constituent of \(\mathcal T\) has determinant weight at most \(2m\beta\), by Proposition proposition-determinant-weights-tensor-operations. Next, Lemma \hyperref[bp66]{14.20} identifies \(H_c^2(U,\mathcal T)\) with \(T_{\pi_1(U)}(-1)\), where \(T\) is the generic fibre of \(\mathcal T\). The geometric fundamental group acts trivially on \(T_{\pi_1(U)}\), so its simple Weil factors are one-dimensional, as one sees by triangularizing the operator induced by a degree-one lift. Let \(b\) be the eigenvalue on such a factor. Since \(T_{\pi_1(U)}\) is a Weil quotient of \(T\), the factor is a constituent of \(\mathcal T\), and its determinant weight is \(\lambda_\iota(b)\). From \(\lambda_\iota(b)=2\log_q|\iota(b)|\le2m\beta\) we get \(|\iota(b)|\le q^{m\beta}\). The twist \((-1)\) multiplies eigenvalues by \(q\), and the claim follows.

\emph{The local factors.} The bound on \(H^2_c\) makes the \(L\)-function holomorphic on a disc, and positivity passes this to each local factor. As \(H_c^0=0\) by Lemma \hyperref[bp66]{14.20}, Theorem \bpout{12.30} gives

\[
L(U_0,\mathcal T,t)
=\prod_{x\in|U_0|}\det(1-\operatorname{Frob}_x t^{\deg(x)}\mid\mathcal T_{\bar x})^{-1}
=\frac{\det(1-Ft\mid H_c^1(U,\mathcal T))}
{\det(1-Ft\mid H_c^2(U,\mathcal T))}.
\]

This is an identity of formal power series, and the product is coefficientwise locally finite because there are only finitely many closed points of bounded degree. Now put \(R_m=q^{-m\beta-1}\). By the bound on \(H_c^2\), the denominator has no zero on the disc \(|t|<R_m\) after applying \(\iota\), so the rational function, and with it its Taylor series, is holomorphic on this disc.

Each local factor \(\det(1-\operatorname{Frob}_x t^{\deg(x)}\mid\mathcal T_{\bar x})\) has no zero on the same disc. Indeed, after applying \(\iota\) the local factors of \(L(U_0,\mathcal T,t)\) are power series with nonnegative coefficients and constant coefficient one by Lemma \hyperref[bp71]{14.22}, so that Lemma \hyperref[bp75]{14.24} applies.

Now let \(\alpha\) be a local eigenvalue of \(\mathcal G\). Then \(\alpha^{2m}\) is an eigenvalue on \(\mathcal T\) by Lemma \bpout{13.5}, and \(1-\iota(\alpha)^{2m}t^{\deg(x)}\) divides the local factor. Its roots have absolute value \(|\iota(\alpha)|^{-2m/\deg(x)}\), so by the previous paragraph this number must be at least \(R_m\), that is,

\[
|\iota(\alpha)|^2\le q_x^{\beta+1/m}.
\]

Letting \(m\) tend to infinity proves the bound.
\end{proof}


% object: sec-Exterior-powers-and-purity (blueprint/chap-Purity-of-real-sheaves/sec-Exterior-powers-and-purity/_section.md), content hidden
\subsection*{\texorpdfstring{14.4\quad Exterior powers and purity}{14.4 Exterior powers and purity}}
\phantomsection\label{bp80}%
\addcontentsline{toc}{subsection}{\texorpdfstring{14.4\quad Exterior powers and purity}{14.4 Exterior powers and purity}}
% text: prose-14-28
\phantomsection\label{bp81}%
Proposition \hyperref[bp78]{14.25} only bounds the local eigenvalues by the largest determinant weight. To bound the eigenvalues of each constituent by its own determinant weight, we apply the proposition to exterior powers, which separate the determinant weights of the simple factors.

Tensor powers cannot do this. For example, suppose \(\mathcal G=\mathcal G_1\oplus\mathcal G_2\) with determinant weights \(\gamma_1>\gamma_2\). In \(\mathcal G^{\otimes m}\) the largest determinant weight \(m\gamma_1\) belongs to the factors of \(\mathcal G_1^{\otimes m}\), while every factor involving \(\mathcal G_2\) lies strictly below it. So the bound of Proposition \hyperref[bp78]{14.25} says nothing sharp about \(\mathcal G_2\).

An exterior power, by contrast, uses each basis vector of \(\mathcal G_1\) at most once. Thus if \(r_1\) is the rank of \(\mathcal G_1\), every graded piece of \(\bigwedge^{r_1+1}\mathcal G\) involves \(\mathcal G_2\), and the summand \(\det\mathcal G_1\otimes\mathcal G_2\) has the largest determinant weight \(r_1\gamma_1+\gamma_2\). Since Corollary \hyperref[bp27]{14.8} gives the eigenvalues of \(\det\mathcal G_1\), a bound for this summand is a bound for \(\mathcal G_2\). We argue on an affine curve, where the proposition applies, and then pass to an arbitrary curve by restricting to affine opens, which preserves irreducibility.

% object: theorem-purity-real-sheaves-curves (blueprint/chap-Purity-of-real-sheaves/sec-Exterior-powers-and-purity/theorem-purity-real-sheaves-curves/_section.md)
\phantomsection\label{bp82}%
\begin{bpstatementit}{Theorem~14.26}{ (Purity of real sheaves on curves)}
On a smooth geometrically connected curve over \(\mathbf F_q\), every irreducible constituent of an \(\iota\)-real lisse Weil sheaf is punctually \(\iota\)-pure of weight equal to its determinant weight. In particular, every \(\iota\)-real lisse Weil sheaf on the curve is mixed.
\end{bpstatementit}

% proof: proof/theorem-purity-real-sheaves-curves
\begin{proof}[{Proof}]\phantomsection\label{bp83}%
We group the constituents by determinant weight and place each group, tensored with the determinants of the groups of larger weight, in an exterior power where it carries the largest determinant weight. There Proposition \hyperref[bp78]{14.25} bounds the eigenvalues of the group by its own determinant weight, and the same argument applied to the dual gives the reverse bound.

Let \(\mathcal G\) be an \(\iota\)-real lisse Weil sheaf on the curve. We first prove the assertion when the curve is affine. There, replacing \(\mathcal G\) by its semisimplification changes neither its local characteristic polynomials nor its irreducible constituents, and it keeps the sheaf real, so we may assume that \(\mathcal G\) is semisimple and write

\[
\mathcal G=\bigoplus_{i=1}^t\mathcal G_i,
\qquad \gamma_1>\cdots>\gamma_t,
\qquad r_i=\operatorname{rank}\mathcal G_i,
\]

where \(\mathcal G_i\) is the sum of the irreducible constituents with determinant weight \(\gamma_i\).

Fix \(i\), and let us find the exterior power in which \(\mathcal G_i\), together with the determinants of \(\mathcal G_1,\dots,\mathcal G_{i-1}\), carries the largest determinant weight. To this end set \(m=\sum_{j<i}r_j\), and put \(B=\sum_{j<i}r_j\gamma_j\).  By Proposition proposition-determinant-weights-tensor-operations, the largest determinant weight of a constituent of \(\bigwedge^{m+1}\mathcal G\) is \(B+\gamma_i\).  Indeed, it is obtained by taking \(n_j=r_j\) for every \(j<i\) and \(n_i=1\), and every other choice gives a smaller sum, because the \(\gamma_j\) are strictly decreasing.  Moreover, exterior powers of a real sheaf are real. Indeed, their eigenvalues are the products of eigenvalues with distinct indices, and complex conjugation permutes this multiset.  Hence Proposition \hyperref[bp78]{14.25} applies to \(\bigwedge^{m+1}\mathcal G\).

The eigenvalues of \(\mathcal G_i\) appear in \(\bigwedge^{m+1}\mathcal G\) not by themselves but multiplied by the eigenvalues of the determinants of the groups of larger weight, and those eigenvalues are known. To make this precise, let \(x\) be a closed point of the curve and \(\alpha\) an eigenvalue of \(\operatorname{Frob}_x\) on \((\mathcal G_i)_{\bar x}\). The decomposition of the exterior power of a direct sum shows that \(\bigwedge^{m+1}\mathcal G\) contains the summand

\[
(\det\mathcal G_1\otimes\cdots\otimes\det\mathcal G_{i-1})
\otimes\mathcal G_i.
\]

If \(\beta_j\) is the eigenvalue of \(\operatorname{Frob}_x\) on \((\det\mathcal G_j)_{\bar x}\), then the product \(\alpha\beta_1\cdots\beta_{i-1}\) is an eigenvalue of \(\operatorname{Frob}_x\) on the stalk of this summand at \(\bar x\). Since \(\det\mathcal G_j\) is pure of weight \(r_j\gamma_j\), we have \(|\iota(\beta_1\cdots\beta_{i-1})|^2=q_x^B\). Now Proposition \hyperref[bp78]{14.25} gives

\[
|\iota(\alpha)|^2q_x^B\le q_x^{B+\gamma_i},
\qquad |\iota(\alpha)|^2\le q_x^{\gamma_i}.
\]

For the reverse bound, note that the dual \(\mathcal G^\vee\) is real as well, with determinant weights \(-\gamma_t>\cdots>-\gamma_1\). Applying the same argument to the summand \(\mathcal G_i^\vee\), we get \(|\iota(\alpha^{-1})|^2\le q_x^{-\gamma_i}\), and hence \(|\iota(\alpha)|^2=q_x^{\gamma_i}\). Since every irreducible constituent of \(\mathcal G\) is a summand of exactly one \(\mathcal G_i\), this proves the theorem on an affine curve.

On an arbitrary curve we restrict to an affine neighbourhood of a given closed point \(x\), which suffices because purity is a condition at each closed point. It remains to check that the restriction changes neither the constituents nor their determinant weights. The neighbourhood is a nonempty open subscheme of a smooth geometrically integral curve, so it is geometrically connected, and restriction to it preserves irreducibility. Indeed, a connected finite étale cover of a normal connected curve is integral, and hence remains connected after restriction to a dense open. So the maps of geometric and of arithmetic fundamental groups are surjective, and since restriction preserves the degree map, the map of Weil groups is surjective too. Hence a subspace of the fibre is stable on the open if and only if it is stable on the curve.

The restriction of a real sheaf is real, and by Corollary \hyperref[bp23]{14.6} the determinant keeps its weight. The affine case therefore applies at \(x\), and since \(x\) was arbitrary, we obtain purity on the whole curve. Finally, the sheaf is mixed because a composition series is a finite filtration with pure quotients.
\end{proof}



\clearpage\phantomsection\hypertarget{bp-outside}{}%
\section*{Outside this export}
\addcontentsline{toc}{section}{Outside this export}
References marked \ensuremath{{}^\ast} are to parts of the blueprint this document leaves out:
\begin{description}
\item[Corollary 2.29] Continuity of cohomology, in Ⅰ The six operations (\bpslug{corollary-{}\allowbreak{}continuity-{}\allowbreak{}cohomology})
\item[Lemma 2.51] Comparison with completed coefficients, in Ⅰ The six operations (\bpslug{lemma-{}\allowbreak{}comparison-{}\allowbreak{}completed-{}\allowbreak{}coefficients})
\item[Proposition 2.55] Continuous cohomology as pro-étale cohomology, in Ⅰ The six operations (\bpslug{proposition-{}\allowbreak{}continuous-{}\allowbreak{}cohomology-{}\allowbreak{}pro-{}\allowbreak{}etale-{}\allowbreak{}cohomology})
\item[Theorem 10.12] Poincaré duality, in Ⅱ Duality and the trace formula (\bpslug{theorem-{}\allowbreak{}poincare-{}\allowbreak{}duality})
\item[Corollary 10.13] Degreewise duality over fields, in Ⅱ Duality and the trace formula (\bpslug{corollary-{}\allowbreak{}degreewise-{}\allowbreak{}duality-{}\allowbreak{}over-{}\allowbreak{}fields})
\item[Theorem 12.30] The trace formula as an Euler product, in Ⅱ Duality and the trace formula (\bpslug{theorem-{}\allowbreak{}trace-{}\allowbreak{}formula-{}\allowbreak{}euler-{}\allowbreak{}product})
\item[Example 13.4] The Kummer extension, in Ⅲ Weights and the Riemann hypothesis (\bpslug{example-{}\allowbreak{}kummer-{}\allowbreak{}extension})
\item[Lemma 13.5] Filtrations and spectra, in Ⅲ Weights and the Riemann hypothesis (\bpslug{lemma-{}\allowbreak{}filtrations-{}\allowbreak{}spectra})
\item[Proposition 13.7] Weights under the usual operations, in Ⅲ Weights and the Riemann hypothesis (\bpslug{proposition-{}\allowbreak{}weights-{}\allowbreak{}under-{}\allowbreak{}usual-{}\allowbreak{}operations})
\item[Chapter 15] Pure sheaves on curves, in Ⅲ Weights and the Riemann hypothesis (\bpslug{chap-{}\allowbreak{}Pure-{}\allowbreak{}sheaves-{}\allowbreak{}on-{}\allowbreak{}curves})
\item[Proposition 15.9] Twisting to weight zero and adding the dual, in Ⅲ Weights and the Riemann hypothesis (\bpslug{proposition-{}\allowbreak{}twisting-{}\allowbreak{}weight-{}\allowbreak{}zero-{}\allowbreak{}adding-{}\allowbreak{}dual})
\item[Chapter 17] The Deligne–Laumon theorem for curves, in Ⅲ Weights and the Riemann hypothesis (\bpslug{chap-{}\allowbreak{}The-{}\allowbreak{}Deligne-{}\allowbreak{}Laumon-{}\allowbreak{}theorem-{}\allowbreak{}for-{}\allowbreak{}curves})
\item[Chapter 18] The Deligne–Laumon theorem in general, in Ⅲ Weights and the Riemann hypothesis (\bpslug{chap-{}\allowbreak{}The-{}\allowbreak{}Deligne-{}\allowbreak{}Laumon-{}\allowbreak{}theorem-{}\allowbreak{}in-{}\allowbreak{}general})
\item[Chapter 19] Proof of the Weil conjectures, in Ⅲ Weights and the Riemann hypothesis (\bpslug{chap-{}\allowbreak{}Proof-{}\allowbreak{}of-{}\allowbreak{}the-{}\allowbreak{}Weil-{}\allowbreak{}conjectures})
\item[Theorem A.2] Nagata's compactification theorem, in Appendices (\bpslug{theorem-{}\allowbreak{}nagata-{}\allowbreak{}s-{}\allowbreak{}compactification-{}\allowbreak{}theorem})
\item[Proposition B.14] The integral Tate module, in Appendices (\bpslug{proposition-{}\allowbreak{}integral-{}\allowbreak{}tate-{}\allowbreak{}module})
\item[Proposition B.15] No Frobenius-fixed vector, in Appendices (\bpslug{proposition-{}\allowbreak{}no-{}\allowbreak{}frobenius-{}\allowbreak{}fixed-{}\allowbreak{}vector})
\item[Appendix C] The Picard scheme of a curve, in Appendices (\bpslug{chap-{}\allowbreak{}The-{}\allowbreak{}Picard-{}\allowbreak{}scheme-{}\allowbreak{}of-{}\allowbreak{}a-{}\allowbreak{}curve})
\item[Corollary C.15] Torsion of the Jacobian and field extension, in Appendices (\bpslug{corollary-{}\allowbreak{}torsion-{}\allowbreak{}jacobian-{}\allowbreak{}field-{}\allowbreak{}extension})
\item[Proposition C.17] $H^1$ of a curve and the Tate module of its Jacobian, in Appendices (\bpslug{proposition-{}\allowbreak{}h-{}\allowbreak{}1-{}\allowbreak{}curve-{}\allowbreak{}tate-{}\allowbreak{}module})
\item[Appendix D] Cohomology of torsion sheaves on curves, in Appendices (\bpslug{chap-{}\allowbreak{}Cohomology-{}\allowbreak{}of-{}\allowbreak{}torsion-{}\allowbreak{}sheaves-{}\allowbreak{}on-{}\allowbreak{}curves})
\item[Proposition D.2] Smooth proper curves, torsion prime to the characteristic, in Appendices (\bpslug{proposition-{}\allowbreak{}smooth-{}\allowbreak{}proper-{}\allowbreak{}curves-{}\allowbreak{}torsion-{}\allowbreak{}prime})
\item[Appendix F] Algebraic groups in characteristic zero, in Appendices (\bpslug{chap-{}\allowbreak{}Algebraic-{}\allowbreak{}groups-{}\allowbreak{}in-{}\allowbreak{}characteristic-{}\allowbreak{}zero})
\item[Lemma F.2] Connected components and images of homomorphisms, in Appendices (\bpslug{lemma-{}\allowbreak{}connected-{}\allowbreak{}components-{}\allowbreak{}images-{}\allowbreak{}homomorphisms})
\item[Lemma F.3] Stable subspaces and the center, in Appendices (\bpslug{lemma-{}\allowbreak{}stable-{}\allowbreak{}subspaces-{}\allowbreak{}center})
\item[Proposition F.8] Lie algebras with a faithful semisimple representation, in Appendices (\bpslug{proposition-{}\allowbreak{}lie-{}\allowbreak{}algebras-{}\allowbreak{}faithful-{}\allowbreak{}semisimple-{}\allowbreak{}representation})
\end{description}
\clearpage\phantomsection%
\section*{References}
\addcontentsline{toc}{section}{References}
\bpbibitem{BBD82}{bp85}{A. Beilinson, J. Bernstein, and P. Deligne. \emph{Faisceaux pervers}, volume 100 of \emph{Astérisque}. 1982.}

\bpbibitem{BS15}{bp86}{B. Bhatt and P. Scholze. The pro-étale topology for schemes. \emph{Astérisque}, 369:99–201, 2015. \href{https://people.mpim-bonn.mpg.de/scholze/proetale.pdf}{Author's manuscript}.}

\bpbibitem{Con07}{bp87}{B. Conrad. Deligne's notes on Nagata compactifications. \emph{Journal of the Ramanujan Mathematical Society}, 22:205–257, 2007. \href{https://math.stanford.edu/\~conrad/papers/nagatafinal.pdf}{Author's manuscript}.}

\bpbibitem{Del74}{bp84}{P. Deligne. La conjecture de Weil. I. \emph{Publ. Math. IHÉS}, 43:273–307, 1974.}

\bpbibitem{Del80}{bp88}{P. Deligne. La conjecture de Weil. II. \emph{Publ. Math. IHÉS}, 52:137–252, 1980.}

\bpbibitem{Den22}{bp89}{C. Deninger. There is no “Weil-”cohomology theory with real coefficients for arithmetic curves. \emph{Ann. Sc. Norm. Super. Pisa Cl. Sci.}, 2022. To appear. arXiv:\href{https://arxiv.org/abs/2204.02714}{2204.02714}.}

\bpbibitem{Deu41}{bp90}{M. Deuring. Die Typen der Multiplikatorenringe elliptischer Funktionenkörper. \emph{Abh. Math. Sem. Univ. Hamburg}, 14:197–272, 1941.}

\bpbibitem{Eke90}{bp91}{T. Ekedahl. On the adic formalism. In \emph{The Grothendieck Festschrift, Vol. II}, volume 87 of \emph{Progress in Mathematics}, pages 197–218. Birkhäuser, 1990.}

\bpbibitem{Has36}{bp92}{H. Hasse. Zur Theorie der abstrakten elliptischen Funktionenkörper. I, II, III. \emph{J. Reine Angew. Math.}, 175:55–62, 69–88, 193–208, 1936.}

\bpbibitem{Hum72}{bp93}{J. E. Humphreys. \emph{Introduction to Lie Algebras and Representation Theory}, volume 9 of \emph{Graduate Texts in Mathematics}. Springer, 1972.}

\bpbibitem{Ill94}{bp94}{L. Illusie. Autour du théorème de monodromie locale. \emph{Astérisque}, 223:9–57, 1994. §4. \href{https://www.numdam.org/item/AST_1994__223__9_0/}{Original article}.}

\bpbibitem{Kat88}{bp95}{N. M. Katz. Travaux de Laumon. \emph{Astérisque}, 161–162:105–132, 1988. Séminaire Bourbaki, exp. 691.}

\bpbibitem{KW01}{bp96}{R. Kiehl and R. Weissauer. \emph{Weil conjectures, perverse sheaves and \(\ell\)-adic Fourier transform}, volume 42 of \emph{Ergebnisse der Mathematik}. Springer, 2001.}

\bpbibitem{Lau87}{bp97}{G. Laumon. Transformation de Fourier, constantes d'équations fonctionnelles et conjecture de Weil. \emph{Publ. Math. IHÉS}, 65:131–210, 1987.}

\bpbibitem{Mil86a}{bp98}{J. S. Milne. \href{https://www.jmilne.org/math/xnotes/AVs.pdf}{Abelian Varieties}. In \emph{Arithmetic Geometry}. 1986. Corrected online version.}

\bpbibitem{Mil86b}{bp99}{J. S. Milne. \href{https://www.jmilne.org/math/xnotes/JVs.pdf}{Jacobian Varieties}. In \emph{Arithmetic Geometry}. 1986. Corrected online version.}

\bpbibitem{Ran39}{bp100}{R. A. Rankin. Contributions to the theory of Ramanujan's function \(\tau(n)\) and similar arithmetical functions. II. \emph{Proc. Cambridge Philos. Soc.}, 35:357–372, 1939.}

\bpbibitem{SGA4}{bp101}{M. Artin, A. Grothendieck, and J.-L. Verdier, editors. \emph{Théorie des topos et cohomologie étale des schémas}, volumes 269, 270, 305 of \emph{Lecture Notes in Mathematics}. Springer, 1972–1973. SGA 4.}

\bpbibitem{SGA4½}{bp102}{P. Deligne. \emph{Cohomologie étale}, volume 569 of \emph{Lecture Notes in Mathematics}. Springer, 1977. SGA 4½; \emph{Théorèmes de finitude}, §3, and \emph{Rapport sur la formule des traces}, §3, and \emph{Sommes trigonométriques}, §1.}

\bpbibitem{SGA5}{bp103}{A. Grothendieck, editor. \emph{Cohomologie \(\ell\)-adique et fonctions \(L\)}, volume 589 of \emph{Lecture Notes in Mathematics}. Springer, 1977. SGA 5.}

\bpbibitem{SGA7II}{bp104}{P. Deligne and N. Katz, editors. \emph{Groupes de monodromie en géométrie algébrique}, volume 340 of \emph{Lecture Notes in Mathematics}. Springer, 1973. SGA 7 II; Exposé XIII, (2.1.7).}

\bpbibitem{Stacks}{bp105}{The Stacks Project Authors. \href{https://stacks.math.columbia.edu}{The Stacks Project}. Individual results are cited by their stable tags in the text.}

\bpbibitem{Var05}{bp106}{Y. Varshavsky. Lefschetz–Verdier trace formula and a generalization of a theorem of Fujiwara. 2005. Version 2. arXiv:\href{https://arxiv.org/abs/math/0505564}{math/0505564}. \href{https://arxiv.org/abs/math/0505564v2}{Author's preprint}.}

\bpbibitem{Wei48a}{bp107}{A. Weil. \emph{Sur les courbes algébriques et les variétés qui s'en déduisent}. Hermann, Paris, 1948.}

\bpbibitem{Wei48b}{bp108}{A. Weil. \emph{Variétés abéliennes et courbes algébriques}. Hermann, Paris, 1948.}

\bpbibitem{Wei49}{bp109}{A. Weil. Numbers of solutions of equations in finite fields. \emph{Bull. Amer. Math. Soc.}, 55:497–508, 1949.}

\end{document}
